% output=pdftex interface=nl 

\environment man-init.tex \hoofdtaal[de]

% For Tobias: 

\let\normalcommaskip=\,

\def\,{\ifmmode\normalcommaskip\else\kern.1em\ignorespaces\fi}

% This is used for abbriviations like d.\,h.

\starttekst

\titelblad
  {PPCH\TeX}
  {ein Makropaket f"ur den Satz\\
   chemischer Strukturformeln}
  {J. Hagen \& T. Burnus\\
   Pragma ADE, Hasselt NL}

\titel{Inhalt}

\stelpaginanummerin[status=stop] \plaatsinhoud \pagina

\titel{Vorwort}

Das Makropaket \PPCHTEX\ wird f"ur den Satz chemischer Strukturformeln 
benutzt. Die Makros basieren auf \PICTEX, ein von Michael Wichura in 
public||domain gebrachtes Makropaket, mit dem Grafiken und Liniendiagramme
erstellt werden k"onnen. Obwohl die erste Version von \PPCHTEX\ nur f"ur
\PICTEX\ entwickelt wurde, kann inzwischen auch \PSTRICKS\ von Timothy Van 
Zandt benutzt werden (wenn auch mit einigen Einschr"ankungen). Man mu"s
aber zugeben, da"s das Ergebnis von \PICTEX\ sch"oner ist.

Die Makros k"onnen in verschiedenen \TEX||Umgebungen verwendet werden; 
dabei wird auf einige \CONTEXT||Module zur"uckgegriffen. Die Makros wurden
so erstellt, da"s eine Erweiterung leicht m"oglich ist. F"ur die 
Interaktivit"at wird auf innerhalb von \CONTEXT\ verwendete Mechanismen
zur"uckgegriffen.

\PPCHTEX\ wurde in erster Linie zum Satz chemischer Strukturformeln 
entworfen, zus"atz\-lich k"onnen aber auch Reaktionsmechanismen wiedergegeben
werden. Die chemischen Strukturformeln k"onnen unterschiedlich dargestellt 
werden, wobei zusammenh"angende Formeln optisch zusammengestellt werden 
k"onnen. Oft benutzte Strukturen k"onnen vordefiniert und abgerufen werden.

Bei der Entwicklung der Makros war Flexiblilit"at, Simplizit"at und
Qualit"at wichtiger als Geschwindigkeit. Die M"oglichkeit Grafiken in
Dateien zu speichern, die \PICTEX\ bietet, wird nicht benutzt. Es hat
sich n"amlich herausgestellt, da"s dies nicht zu einem
Geschwindigkeitsvorteil f"uhrt.

Die erste Version von \PPCHTEX\ erschien 1995. Diese Anleitung beschreibt die
zweite Version, die dank der vielen Vorschl"age von Tobias Burnus, Dirk 
Kuypers und Ton Otten die doppelte Funktionalit"at hat.

\pagina \stelpaginanummerin[status=start]

\vervolgblad
  {Grundlagen}
  []
  [SIX,B,C,R1356,RZ1356]
  [NO_2,NO_2,O_2N,CH_3]

\hoofdstuk{Strukturen}

Die Anzahl der Befehle, die zum Satz von Strukturformeln gebraucht werden,
beschr"anken sich auf vier.\voetnoot{In dieser Anleitung 
bezieht sich der Begriff Struktur nur auf chemische Strukturen und nicht auf
die Struktur des Textes, in der die Formel plaziert wird.} Im folgenden 
Beispiel werden all diese Befehle benutzt:

\startbuffer
\stelchemiein[assenstelsel=aan,kader=aan]

\startchemie
  \chemie[SIX,B,R,RZ][1,2,3,4,5,6]
\stopchemie
\stopbuffer

\startfiguurtekst
  [links][]
  {geen}
  {\haalbuffer}
\voorbeeld[vb:zesring]
\starttypen
\stellechemieein[achsen=an,rahmen=an]
\startchemie
  \chemie[SIX,B,R,RZ][1,2,3,4,5,6]
\stopchemie
\stoptypen
\stopfiguurtekst

Mit \type{\stellechemieein} k"onnen verschiedene Optionen des Setzsystems
eingestellt werden. Die mit diesem Befehl vorgenommen Einstellungen gelten
f"ur alle folgenden Formeln.\voetnoot{Nat"urlich kann die G"ultigkeit 
durch \type{{ }} und den Gruppierungsmakros \type{..group} begrenzt werden.
Die Einstellungen k"onnen auch direkt nach \type{startchemie} gegeben 
werden; sie gelten dann nur f"ur die jeweilige Formel.} 

\startbuffer
\startchemie[kader=aan,breedte=passend,hoogte=passend]
  \chemie[CARBON,CB1][A,B,C,D,E,F]
\stopchemie
\stopbuffer

\startfiguurtekst
  [links][]
  {geen}
  {\haalbuffer}
\voorbeeld
\starttypen
\startchemie[rahmen=an,breite=passend,hoehe=passend]
  \chemie[CARBON,CB1][A,B,C,D,E,F]
\stopchemie
\stoptypen
\stopfiguurtekst

Wie man an beiden Beispielen erkennt ist \type{\chemie} der zentrale Befehl.
Dieser Befehl, der mehrmals innerhalb eines
\type{\start}||\type{\stop}||Paars vorkommen kann, hat ein oder zwei
Argumente. Diese werden zwischen~\type{[ ]} angegeben. Das erste Argument
bezieht sich auf die zu zeichnenden Bindungen, wogegen das zweite die Atome
oder Molek"ule, die dargestellt werden sollen, enth"alt. Der Text wird im
mathematischen Modus gesetzt, somit kann alles, was zwischen~\type{$ $}
erlaubt ist, verwendet werden. 

Gehen wir nun das erste Beispiel durch. Zuerst wird das Schl"usselword
\type{SIX} angegeben. Durch Verwendung dieses Wortes teilen wir dem Makro mit,
da"s wir eine Sechsringstruktur zeichnen wollen. Ebenso k"onnen wir
\type{ONE}, \type{THREE}, \type{FOUR}, \type{FIVE}, \type{EIGHT}, 
\type{CARBON}, \type{NEWMAN} und \type{CHAIR} verwenden.

\def\structuurtje[#1]%
  {\vbox
     {\startchemie[kader=aan,schaal=klein,formaat=klein]%
      \chemie[#1][1,2,3,4,5,6,7,8]%
      \stopchemie}}

\plaatsfiguur{geen}
 \startcombinatie[4*3]
   {\structuurtje[ONE,SB,Z]}           {One}
   {\structuurtje[THREE,B,R,RZ]}       {Three}
   {\structuurtje[FOUR,B,R,RZ]}        {Four}
   {\structuurtje[FIVE,B,R,CRZ]}       {Five}
   {\structuurtje[SIX,B,R,RZ]}         {Six}
   {\structuurtje[EIGHT,B]}            {Eight}
   {\structuurtje[FIVE,FRONT,B,-R,+R]} {Five Front}
   {\structuurtje[SIX,FRONT,B,-R,+R]}  {Six Front}
   {\structuurtje[CARBON,CB]}          {Carbon}
   {\structuurtje[NEWMAN,STAGGER,CB]}  {Newman Stacker}
   {\structuurtje[NEWMAN,ECLIPSE,CB]}  {Newman Eclipse}
   {\structuurtje[CHAIR,B,-R,+R]}      {Chair}
 \stopcombinatie

Bez"uglich der Abmessungen f"allt \type{CHAIR} auf. Auch in anderen 
Bereichen unterscheidet sich diese Struktur von den anderen. So ist 
beispielsweise das Drehen und Kombinieren nicht m"oglich.

\startbuffer
\startchemie[kader=aan,schaal=klein]
  \chemie
    [CHAIR,B,+R,-R,
     +RZ1,+RZ2,+RZ3,+RZ4,+RZ5,+RZ6,
     -RZ1,-RZ2,-RZ3,-RZ4,-RZ5,-RZ6]
    [a,b,c,d,e,f,1,2,3,4,5,6]
\stopchemie
\stopbuffer

\startfiguurtekst
  [links][]
  {geen}
  {\haalbuffer}
\voorbeeld
\starttypen
\startchemie[rahmen=an,format=klein]
  \chemie
    [CHAIR,B,+R,-R,
     +RZ1,+RZ2,+RZ3,+RZ4,+RZ5,+RZ6,
     -RZ1,-RZ2,-RZ3,-RZ4,-RZ5,-RZ6]
    [a,b,c,d,e,f,1,2,3,4,5,6]
\stopchemie
\stoptypen
\stopfiguurtekst

Innerhalb dieser Strukturen werden die Bindungen zwischen den \chemie{C}||Atomen auf
eine vergleichbare Weise angegeben. In diesem Beispiel verwenden wir~\type{B}
und~\type{R}. Die Bindungen, die nummeriert sind, k"onnen verschiedenartig
angegeben werden:

\starttypen
\chemie[SIX,B1,B2,B3,B4,B5,B6]
\chemie[SIX,B135]
\chemie[SIX,B1..5]
\stoptypen

Diese Befehle erstellen Teile eines Sechsrings. Mit \type{R} und \type{RZ}
k"onnen Substituenten an den Sechsring plaziert werden. Der Parameter \type{R}
zeichnet eine Bindung von einer Sechs\-ring\-ecke ($\angle~120^\circ$) zu einem 
Substituenten. Die jeweilige Ringecke wird mit einer Nummer angegeben.

\starttypen
\chemie[SIX,B1..6,R1..6]
\stoptypen

Dieser obrige Befehl zeichnet nur die Bindung zum Substituenten. Der 
Substituent selbst wird mit \type{RZ} angegeben. Auch hier mu"s, um die 
Position zu bestimmen, die Nummer angegeben werden. Als zweites optionales
Argument werden die Substituenten angegeben.

\startbuffer
\startchemie[kader=aan,breedte=6000]
  \chemie[SIX,B1..6,R1..6,RZ1..3][CH_3,CH_3,OH]
\stopchemie
\stopbuffer

\startfiguurtekst
  [links][]
  {geen}
  {\haalbuffer}
\voorbeeld
\starttypen
\startchemie[rahmen=an,breite=6000]
  \chemie[SIX,B1..6,R1..6,RZ1..3][CH_3,CH_3,OH]
\stopchemie
\stoptypen
\stopfiguurtekst

Wenn das zweite Argument weggelassen wird, wird kein Text platziert, d.~h.
der Parameter \type{RZ1..3} hat keinen Effekt.

\hoofdstuk{Bindungen}

In diesem Kapitel geben wir eine "Ubersicht "uber die Bindungen, die man
bei unterschiedlichen Strukturen verwenden kann. In den Beispielen und 
"Ubersichten in sp"ateren Kapiteln werden wir n"aher auf die Parameter 
eingehen.

In der linken Spalte stehen die vollst"andigen Bindungen, in der rechten
die verk"urzten. Letztere erm"oglichen es, da"s Atome und Molek"ule an
diese platziert werden k"onnen. Die Bindungen k"onnen auf beiden Seiten,
links (\type{-}) sowie rechts (\type{+}), verk"urzt sein.

\plaatstabel
{Einfachbindungen.}
\starttabel[|rT|lw(4cm)|rT|lw(4cm)|]
\HL
\VL ~~~B \VL Bond          \VL ~~SB \VL Single Bond       \VL\FR
\VL ~~BB \VL Bold Bond     \VL ~-SB \VL Left Single Bond  \VL\MR
\VL ~~HB \VL Hydrogen Bond \VL ~+SB \VL Right Single Bond \VL\LR
\HL
\stoptabel

In diesem Beispiel wurden einige der Bindungen kombiniert:

\startbuffer
\startchemie[kader=aan]
  \chemie[ONE,SD1,SB4,BB2,SB7,Z01247][C,H,H,H,H]
\stopchemie
\stopbuffer

\startfiguurtekst
  [links][]
  {geen}
  {\haalbuffer}
\voorbeeld
\starttypen
\startchemie[rahmen=an]
  \chemie[ONE,SD1,SB4,BB2,SB7,Z01247][C,H,H,H,H]
\stopchemie
\stoptypen
\stopfiguurtekst

Einer Bindung kann eine oder mehrere Zahlen bzw. ein Zahlenbereich folgen,
zum Beispiel: \type{B1}, \type{B135} und \type{B1..5}. Wenn alle Bindungen
n"otig sind, kann einfach nur~\type{B} angegeben werden.

In einem Ring kann eine zus"atzliche Bindung angegeben werden (f"uhrt zu 
einer Doppelbindung); zwischen Atomen und Molek"ulen k"onnen Doppel|| und 
Dreifachbindungen gezeichnet werden.

\plaatstabel
{Mehrfachbindungen.}
\starttabel[|rT|lw(4cm)|rT|lw(4cm)|]
\HL
\VL ~~EB \VL Extra Bond    \VL ~~DB \VL Double Bond       \VL\FR
\VL      \VL               \VL ~~TB \VL Triple Bond       \VL\LR
\HL
\stoptabel

Freie Elektronen k"onnen verschiedenartig angegeben werden. Die 
zugeh"origen Parameter beginnen mit einem \type{E}.

\plaatstabel
{Freie Elektronen.}
\starttabel[|rT|lw(4cm)|rT|lw(4cm)|]
\HL
\VL ~~ES \VL Extra Single  \VL ~~ED \VL Extra Double \VL\FR
\VL ~~EP \VL Extra Pair    \VL ~~ET \VL Extra Triple \VL\LR
\HL
\stoptabel

Das untenstehende Kohlenstoffatom hat nur aus Demonstrationszwecken
8~Au"sen\-elek\-tronen:

\startbuffer
\startchemie[kader=aan]
  \chemie[ONE,Z0,ES1,ED3,ET5,EP7][C]
\stopchemie
\stopbuffer

\startfiguurtekst
  [links][]
  {geen}
  {\haalbuffer}
\voorbeeld
\starttypen
\startchemie[rahmen=an]
  \chemie[ONE,Z0,ES1,ED3,ET5,EP7][C]
\stopchemie
\stoptypen
\stopfiguurtekst

Eine Ringbindung kann \citaat{kurzgeschlossen} werden. Dies kommt vor allem
bei Sechsringen vor. In diesem Fall wird das Atom, das wegfallen soll, 
angegeben. Es es auch m"oglich einen Kreis innerhalb eines Sechsrings zu 
zeichen.

\plaatstabel
{Besondere Bindungen.}
\starttabel[|rT|lw(4cm)|rT|lw(4cm)|]
\HL
\VL ~~SS \VL Short Shortcut       \VL ~~~S \VL Shortcut            \VL\FR
\VL ~-SS \VL Left  Short Shortcut \VL ~MID \VL Mid Shortcut        \VL\MR
\VL ~+SS \VL Right Short Shortcut \VL MIDS \VL Closed Mid Shortcut \VL\LR
\HL
\stoptabel

\plaatstabel
{Kreisf"ormige Bindungen.}
\starttabel[|rT|lw(4cm)|rT|lw(4cm)|]
\HL
\VL ~~~C \VL Circle               \VL ~~CD \VL Dashed Circle          \VL\FR
\VL ~~CC \VL Shifted Circle       \VL ~CCD \VL Dashed Shifted Circle  \VL\LR
\HL
\stoptabel

Obwohl sich die Bedeutung der kreisf"ormigen \type{C..}||Parameter leicht 
erraten la"st, ist hier trotzdem ein Beispiel.

\startbuffer
\startchemie[kader=aan]
  \chemie
    [SIX,B,ER6,CCD12346,Z0,PB:RZ6,ONE,SB8,EP6,Z0,ZT6,Z8,PE]
    [\ominus,O,\oplus,H]
\stopchemie
\stopbuffer

\startfiguurtekst
  [links][]
  {geen}
  {\haalbuffer}
\voorbeeld
\starttypen
\startchemie[rahmen=an]
  \chemie
    [SIX,B,ER6,CCD12346,Z0,PB:RZ6,ONE,SB8,EP6,Z0,ZT6,Z8,PE]
    [\ominus,O,\oplus,H]
\stopchemie
\stopchemie
\stoptypen
\stopfiguurtekst

An allen Ringecken k"onnen Substituenten gebunden werden. Wobei Substituent
im weiteren Sinn zu verstehen ist. Die Bindungen k"onnen, abh"angig von
der Anwesenheit von Atomen und Molek"ulen, kurz oder lang sein. Wir werden
diesen Parametern oft begegnen.

\plaatstabel
{Bindungen zu Substituenten.}
\starttabel[|rT|lw(4cm)|rT|lw(4cm)|]
\HL
\VL ~~~R \VL Radical       \VL ~~SR \VL Single Radical       \VL\FR
\VL ~~-R \VL Left Radical  \VL ~-SR \VL Single Left Radical  \VL\MR
\VL ~~+R \VL Right Radical \VL ~+SR \VL Single Right Radical \VL\LR
\HL
\stoptabel

Neben dieser Bindungen gibt es noch andere Einfachbindungen.

\plaatstabel
{Besondere Bindungen zu Substituenten.}
\starttabel[|rT|lw(4cm)|rT|lw(4cm)|]
\HL
\VL ~~RD \VL Radical Dashed       \VL ~~RB\VL Radical Bold       \VL\FR
\VL ~-RD \VL Left Radical Dashed  \VL ~-RB\VL Left Radical Bold  \VL\MR
\VL ~+RD \VL Right Radical Dashed \VL ~+RB\VL Right Radical Bold \VL\LR
\HL
\stoptabel

Die Radikale k"onnen stets in drei Richtungen gezeichnet werden. Die linke
und rechte Variante ($-$ und $+$) werden aber sehr selten ben"otigt.

\startbuffer
\startchemie[kader=aan]
  \chemie[SIX,B,R14,RD25,RB36]
\stopchemie
\stopbuffer

\startfiguurtekst
  [links][]
  {geen}
  {\haalbuffer}
\voorbeeld
\starttypen
\startchemie[rahmen=an]
  \chemie[SIX,B,R14,RD25,RB36]
\stopchemie
\stoptypen
\stopfiguurtekst


\plaatstabel
{Weitere besondere Bindungen zu Substituenten.}
\starttabel[|rT|lw(4cm)|rT|lw(4cm)|]
\HL
\VL ~~SD \VL Single Dashed \VL ~LDD \VL Double Dashed Left  \VL\FR
\VL ~~OE \VL Open Ended    \VL ~RDD \VL Double Dashed Right \VL\LR
\HL
\stoptabel

Hier folgt ein Beispiel f"ur ein {\em Open End}. Man sieht einen Sechsring
(\type{SIX}), an dem mehrere \type{ONE}s gebunden sind. Auf die Bedeutung 
von \type{PB} gehen wir sp"ater ein.

\startbuffer
\startchemie
    [breedte=5000,boven=2500,onder=1500,kader=aan]
  \chemie
    [SIX,B,C,R6,
     PB:RZ6,ONE,CZ0,OE1,SB5,MOV5,CZ0,OFF5,OE5,PE]
    [CH,CH_2]
\stopchemie
\stopbuffer

\startfiguurtekst
  [links][]
  {geen}
  {\haalbuffer}
\voorbeeld
\starttypen
\startchemie
    [breite=5000,oben=2500,unten=1500,rahmen=an]
  \chemie
    [SIX,B,C,R6,
     PB:RZ6,ONE,CZ0,OE1,SB5,MOV5,CZ0,OFF5,OE5,PE]
    [CH,CH_2]
\stopchemie
\stoptypen
\stopfiguurtekst

Nat"urlich k"onnen die Substituenten auch mit einer Doppelbindung an die
Struktur gebunden werden.

\plaatstabel
{Mehrfachbindungen zu Substituenten.}
\starttabel[|rT|lw(4cm)|rT|lw(4cm)|]
\HL
\VL ~~ER \VL Extra Radical \VL ~~DR \VL Double Radical \VL\SR
\HL
\stoptabel

An die Bindungen kann Text gebunden werden. Dieser steht bei \type{\chemie}
im zweiten Klammerpaar.

\plaatstabel
{Atome und Molek"ule (Radikale).}
\starttabel[|rT|lw(4cm)|rT|lw(4cm)|]
\HL
\VL ~~~Z \VL Atom        \VL ~~RZ \VL Radical Atom       \VL\FR
\VL ~CRZ \VL Center Atom \VL ~-RZ \VL Left Radical Atom  \VL\MR
\VL MIDZ \VL Mid Atom    \VL ~+RZ \VL Right Radical Atom \VL\LR
\HL
\stoptabel

Bei diesen Parametern schlie"st \type{RZ} an \type{R}, das zuerst angegeben
werden mu"s, an. Der Parameter \type{MID} kann nur bei einem Sechsring 
(\type{SIX}) verwendet werden. In dem folgenden Beispiel wird sowohl 
\type{MID} als auch \type{MIDZ} verwendet. Bei diesen Parametern wird keine
Positionzahl angegeben.

\startbuffer
\startchemie[kader=aan]
  \chemie[SIX,B,MID,MIDZ][\SL{CH_2}]
\stopchemie
\stopbuffer

\startfiguurtekst
  [links][]
  {geen}
  {\haalbuffer}
\voorbeeld
\starttypen
\startchemie[rahmen=an]
  \chemie[SIX,B,MID,MIDZ][\SL{CH_2}]
\stopchemie
\stoptypen
\stopfiguurtekst

Die Atome/Molek"ule werden im Uhrzeigersinn nummeriert. Auch hier sind
Kombinationen erlaubt. Die Position \type{0} (null) plaziert den Text in
die Mitte einer Struktur.

Wir k"onnen die Atombindungen nummerieren und beschriften. Dies geschieht
mit \type{ZN} und \type{ZT}:

\plaatstabel
{Beschriftung und Nummerierung der Bindungen.}
\starttabel[|rT|lw(4cm)|rT|lw(4cm)|]
\HL
\VL ~~ZN \VL Atom Number \VL ~~ZT \VL Atom Text \VL\SR
\HL
\stoptabel

Die obrigen Parameter haben bei \type{ONE} noch eine oben (top) und unten 
(bottom) Variante.

\plaatstabel
{Zus"atzliche Beschriftung und Nummerierung der Bindungen.}
\starttabel[|rT|lw(4cm)|rT|lw(4cm)|]
\HL
\VL ~ZTN \VL Atom Top Number    \VL ~ZTT \VL Atom Top Text    \VL\FR
\VL ~ZBN \VL Atom Bottom Number \VL ~ZBT \VL Atom Bottom Text \VL\LR
\HL
\stoptabel

Bei \type{ZTN} und \type{ZBN} werden die Nummern automatisch erstellt. Beim
anderen Parameter wird der "ubergebene Text benutzt.

\startbuffer
\startchemie[kader=aan]
  \chemie[ONE,SB,Z0,ZTT][C,a,b,c,d,e,f,g,h]
\stopchemie
\stopbuffer

\startfiguurtekst
  [links][]
  {geen}
  {\haalbuffer}
\voorbeeld
\starttypen
\startchemie[rahmen=an]
  \chemie[ONE,SB,Z0,ZTT][C,a,b,c,d,e,f,g,h]
\stopchemie
\stoptypen
\stopfiguurtekst

Bei \type{SIX} und \type{FIVE} kann man auch Radikalebindungen nummerieren. 
Dazu benutzen wir \type{RN} und \type{RT}.

\plaatstabel
{Beschriftungen und Nummern.}
\starttabel[|rT|lw(4cm)|rT|lw(4cm)|]
\HL
\VL ~~RN \VL Radical Number        \VL ~~RT \VL Radical Text        \VL\FR
\VL ~RTN \VL Radical Top Number    \VL ~RTT \VL Radical Top Text    \VL\MR
\VL ~RBN \VL Radical Bottom Number \VL ~RBT \VL Radical Bottom Text \VL\LR
\HL
\stoptabel

Zuletzt noch Symbole, die f"ur den Drehsinn einer Struktur verwendet werden.

\plaatstabel
{Anzeige des Drehsinns.}
\starttabel[|rT|lw(4cm)|rT|lw(4cm)|]
\HL
\VL ~~AU \VL Arrow Up \VL ~~AD \VL Arrow Down \VL\SR
\HL
\stoptabel

Die Pfeile werden zwischen die Radikale gezeichnet.

\startbuffer
\startchemie[kader=aan]
  \chemie[SIX,B,AU]
\stopchemie
\stopbuffer

\startfiguurtekst
  [links][]
  {geen}
  {\haalbuffer}
\voorbeeld
\starttypen
\startchemie[rahmen=an]
  \chemie[SIX,B,AU]
\stopchemie
\stoptypen
\stopfiguurtekst

Beim Platzieren der Atome und Molek"ule in einem Textes, wird so gut wie 
m"oglich die Gr"o"se der Struktur bestimmt. Die Breite des \chemie{C} und
die H"ohe von \chemie{C^n_m} spielen dabei eine Rolle. Der Mechanismus kann
aber noch verfeinert werden.

\hoofdstuk{Frontalansicht}

Die Strukturen \type{FIVE} und \type{SIX} sind auch als Frontalansicht
darstellbar (die der {\sc Haworth}||Projektion bei Sacchariden entspricht).
Als Beispiel f"ur die Strukturen werden unten zwei Monosaccheride 
dargestellt.

Die Frontalansichten k"onnen nicht rotiert werden; auch sind Verkn"upfungen
mit anderen Strukturen (noch) nicht m"oglich.

\startbuffer
\startchemie[hoogte=4500,onder=2500]
  \bottext{$\beta$-{\sc d}-Fructopyranose}
  \chemie
    [SIX,FRONT,BB1236,+SB4,-SB5,Z5][O]
  \chemie
    [SIX,FRONT,+R12346,+RZ12346][\SR{HO},H,H,H,OH]
  \chemie
    [SIX,FRONT,-R12346,-RZ12346][H,OH,\SR{HO},H,CH_2OH]
\stopchemie
\stopbuffer

\startfiguurtekst
  [links][]
  {geen}
  {\haalbuffer}
\voorbeeld
\starttypen
\startchemie[hoehe=4500,unten=2500]
  \bottext{$\beta$-{\sc d}-Fructopyranose}
  \chemie
    [SIX,FRONT,BB1236,+SB4,-SB5,Z5][O]
  \chemie
    [SIX,FRONT,+R12346,+RZ12346][\SR{HO},H,H,H,OH]
  \chemie
    [SIX,FRONT,-R12346,-RZ12346][H,OH,\SR{HO},H,CH_2OH]
\stopchemie
\stoptypen
\stopfiguurtekst


Das Positionieren der Radikale ist ein Kombination von
Durchf"uhrbarkeit und Qualit"at. 
Das folgende Beispiel macht dies noch deutlicher:

\startbuffer
\startchemie[hoogte=4500,onder=2500]
  \bottext{$\alpha$-{\sc d}-Fructofuranose}
  \chemie
    [FIVE,FRONT,BB125,+SB3,-SB4,Z4][O]
  \chemie
    [FIVE,FRONT,+R1235,+RZ1235][\SR{HO},H,\SR{HOH_2C},CH_2OH]
  \chemie
    [FIVE,FRONT,-R1235,-RZ1235][OH,H,\SR{HO},H,CH_2OH]
\stopchemie
\stopbuffer

\startfiguurtekst
  [links][]
  {geen}
  {\haalbuffer}
\voorbeeld
\starttypen
\startchemie[hoehe=4500,unten=2500]
  \bottext{$\alpha$-{\sc d}-Fructofuranose}
  \chemie
    [FIVE,FRONT,BB125,+SB3,-SB4,Z4][O]
  \chemie
    [FIVE,FRONT,+R1235,+RZ1235][\SR{HO},H,\SR{HOH_2C},CH_2OH]
  \chemie
    [FIVE,FRONT,-R1235,-RZ1235][OH,H,\SR{HO},H,CH_2OH]
\stopchemie
\stoptypen
\stopfiguurtekst

\hoofdstuk{Definitionen} 

Es ist m"oglich eine Strukturbibliothek zu erstellen. Diese 
Strukturen k"onnen wir dann sp"ater wieder aufrufen und erweitern.
Weiter k"onnen sie zur Blockbildung beim Aufbau komplizierter Strukturen
dienen. Strukturen k"onnen mit dem \TEX-Befehl \type{\def} vordefiniert 
werden.

\starttypen
\def\sechsring{\chemie[SIX,B,R,RZ]}
\stoptypen

An Stelle von \type{\def} kann auch der folgende Befehl benutzt werden.
In diesem Fall wird eine Meldung gegeben, da"s die Definition bereits
besteht.

\starttypen
\definierechemie[sechsring]
  {\chemie[SIX,B,R,RZ]}
\stoptypen

Obwohl beide Definitionen m"oglich sind, ist die zweite robuster, da
diese darauf achtet bestehende Befehle nicht zu "uberschreiben.

Der Befehl \type{\chemie[sechsring]} liefert einen Sechsring ohne 
Substituenten. Es wurde kein zweites Argument angegben.

\startbuffer
\definieerchemie[zesring]
  {\chemie[SIX,B,R,RZ]}

\startchemie[kader=aan,breedte=6000]
  \chemie[zesring]
\stopchemie
\stopbuffer

\startfiguurtekst
  [links][]
  {geen}
  {\haalbuffer}
\voorbeeld
\starttypen
\definierechemie[sechsring]
  {\chemie[SIX,B,R,RZ]}

\startchemie[rahmen=an,breite=6000]
  \chemie[sechsring]
\stopchemie
\stoptypen
\stopfiguurtekst

Um Substituenten hinzuzuf"ugen, m"ussen wir das Folgende machen:

\startbuffer
\definieerchemie[zesring]
  {\chemie[SIX,B,R,RZ]}

\startchemie[kader=aan,breedte=6000]
  \chemie[zesring][R_1,R_2,R_3,R_4,R_5,R_6]
\stopchemie
\stopbuffer

\startfiguurtekst
  [links][]
  {geen}
  {\haalbuffer}
\voorbeeld
\starttypen
\definierechemie[sechsring]
  {\chemie[SIX,B,R,RZ]}

\startchemie[rahmen=an,breite=6000]
  \chemie[sechsring][R_1,R_2,R_3,R_4,R_5,R_6]
\stopchemie
\stoptypen
\stopfiguurtekst

Die Struktur \type{sechsring} kann auch ohne Substituenten (\type{RZ}) 
definiert werden. In diesm Fall werden durch \type{\chemie[sechsring]} keine
Substituenten gezeichnet. Trotzdem k"onnen diese noch gezeichnet werden,
wie es im folgenden Beispiel gezeigt wird.

\startbuffer
\definieerchemie[zesring]
  {\chemie[SIX,B,R]}

\startchemie[kader=aan,breedte=6000]
  \chemie[zesring,RZ][A,B,C,D,E,F]
\stopchemie
\stopbuffer

\startfiguurtekst
  [links][]
  {geen}
  {\haalbuffer}
\voorbeeld
\starttypen
\definierechemie[sechsring]
  {\chemie[SIX,B,R]}

\startchemie[rahmen=an,breite=6000]
  \chemie[sechsring,RZ][A,B,C,D,E,F]
\stopchemie
\stoptypen
\stopfiguurtekst

Im Prinzip sind die M"oglichkeiten unbegrenzt. Man sollte sich aber 
bewu"st sein, da"s die Argumente im zweiten Klammerpaar in der 
Reihenfolge der Parameter der ersten verwendet werden.

In einer Definition k"onnen aber auch Atome und Molek"ule (Text) plaziert
werden.

\starttypen
\definierechemie[sechsring]
  {\chemie[SIX,B,R,RZ135][R_1,R_3,R_5]}
\stoptypen

Hier werden immer drei Substituenten plaziert. Wenn wir mehr Substituenten
platzieren wollen, m"ussen wir explizit angeben, da"s wir eine 
Sechsringstruktur (\type{SIX}) verwenden.

\startbuffer
\definieerchemie[zesring]
  {\chemie[SIX,B,R,RZ135][R_1,R_3,R_5]}

\startchemie[kader=aan,breedte=6000]
  \chemie[zesring,SIX,RZ246][A,B,C]
\stopchemie
\stopbuffer

\startfiguurtekst
  [links][]
  {geen}
  {\haalbuffer}
\voorbeeld
\starttypen
\definierechemie[sechsring]
  {\chemie[SIX,B,R,RZ135][R_1,R_3,R_5]}

\startchemie[rahmen=an,breite=6000]
  \chemie[sechsring,SIX,RZ246][A,B,C]
\stopchemie
\stoptypen
\stopfiguurtekst

In den Definitionen ist \type{\chemie[ ]} global (d.~h. das sich die
Struktur \type{SIX} gemerkt wird) und \type{\chemie[ ][ ]} lokal. Die Idee 
dahinter ist, da"s im ersten Fall eine Folge von Befehlen angegeben wird,
wogegen das zweite als v"ollig unabh"angige Struktur angesehen wird.

In einer Definition kann \type{\chemie} mehrmals vorkommen. Das vorige
Beispiel kann auch so angegeben werden:

\startbuffer
\definieerchemie[zesring]
  {\chemie[SIX,B,R,RZ135][R_1,R_3,R_5]
   \chemie[SIX,RZ246]}

\startchemie[kader=aan,breedte=6000]
  \chemie[zesring][A,B,C]
\stopchemie
\stopbuffer

\startfiguurtekst
  [links][]
  {geen}
  {\haalbuffer}
\voorbeeld
\starttypen
\definierechemie[sechsring]
  {\chemie[SIX,B,R,RZ135][R_1,R_3,R_5]
   \chemie[SIX,RZ246]}

\startchemie[rahmen=an,breite=6000]
  \chemie[sechsring][A,B,C]
\stopchemie
\stoptypen
\stopfiguurtekst

Wenn \TEX\ meldet, da"s es auf einen unbekannten Befehl gestossen ist,
dann hat man wahrscheinlich vergessen \type{SIX}, \type{FIVE} oder ein 
vergleichbaren Strukturparameter anzugeben.

\hoofdstuk{Kombinieren}

Strukturen k"onnen zu komplexen Verbindungen verbunden werden. Das 
Verschieben einer Struktur geschieht mit \type{MOV}, \type{ROT}, 
\type{ADJ} und \type{SUB}.

\plaatstabel
{Verschieben und Drehen.}
\starttabel[ | r T | l w(2cm) | p(8cm) | ]
\HL
\VL MOV \VL Move       \VL Verschieben der selben Art von Struktur in Richtung
                           einer Bindung\VL\FR 
\VL ADJ \VL Adjace     \VL Verschieben einer anders artigen Struktur, die 
                           \"uber eine Bindung verbunden wird, in 
                           Richtung der $x$- oder $y$-Achsen\VL\MR
\VL SUB \VL Substitute \VL Verschieben einer Struktur relativ zu einer anderen
                           in Richtung der $x$- oder $y$-Achse\VL\MR 
\VL ROT \VL Rotate     \VL Drehen einer Struktur\VL\LR 
\HL
\stoptabel

Die obrigen vier Parameter haben innerhalb der verschieden Strukturen 
unterschiedliche Wirkung. So unterscheidet sich beispielsweise der
Drehwinkel bei \type{\chemie[FIVE,ROT1,B]} von dem bei 
\type{\chemie[SIX,ROT1,B]}.

\type{ONE} ist auf die Parameter \type{MOV} und \type{DIR} beschr"ankt.
Beide Befehle haben den selben Effekt mit Au"snahme der L"ange der
Verschiebung. Kleine Verschiebungen sind mit \type{OFF} m"oglich.

\plaatstabel
{Verschieben.}
\starttabel[ | r T | l w(2cm) | p(8cm) | ]
\HL
\VL DIR \VL Direction \VL Wie \type{MOV}, aber zum diagonalen Verschieben von
                          Strukturen\VL\FR
\VL OFF \VL Offset    \VL Kleine Verschiebungen eines Atom innerhalb eines 
                          Molek\"uls\VL\LR 
\HL
\stoptabel

Innerhalb von \type{CARBON} ist es auch m"oglich die Struktur zu spiegeln.
Dies geschieht mit \type{MIR}.

\plaatstabel
{Spiegelen.}
\starttabel[ | r T | l w(2cm) | p(8cm) | ]
\HL
\VL MIR \VL Mirror     \VL Spiegeln einer Struktur\VL\SR
\HL
\stoptabel

Mit einer Ziffer wird die Richtung der Verschiebung bzw. das Ma"s der
Drehung angegeben. Da diese Parameter stark von der Struktur abh"angen,
m"ussen sie angegeben werden, bevor Bindungen und Text gezeichnet wurden.
Es macht einen Unterschied, ob \type{\chemie[FIVE,B,ROT1,R]} oder 
\type{\chemie[FIVE,ROT1,B,R]} angegeben wird. Das erste liefert ein
unerw"unschtes Resultat.

\startbuffer
\startchemie[kader=aan,breedte=4000,rechts=3000]
  \chemie[SIX,B,MOV1,B]
\stopchemie
\stopbuffer

\startfiguurtekst
  [links][]
  {geen}
  {\haalbuffer}
\voorbeeld
\starttypen
\startchemie[rahmen=an,breite=4000,rechts=3000]
  \chemie[SIX,B,MOV1,B]
\stopchemie
\stoptypen
\stopfiguurtekst

Zuerst wurde ein Sechsring gezeichnet: \type{SIX,B}, die aktuelle Position
wurde dann in Richtung der Bindung~1 verschoben: \type{MOV1} und ein zweiter 
Ring gezeichnet:~\type{B}. Durch die Verschiebung mit \type{MOV} kann der zweite
Sechsring an allen sechs Kanten plaziert werden, wogegen \type{ADJ} nur
in Achsenrichtung ($x$, $-x$, $y$, $-y$) verschieben kann.
Bei einem Sechsring fallen die Verschiebungen sogar zusammen, das obenstehende
Beispiel h"atte auch so erreicht werden k"onnen: \type{[SIX,B,ADJ1,B]}.

Auch k"onnen verschiedene Strukturen kombiniert werden. An einer Struktur
\type{SIX} kann beispielsweise \type{FIVE} gekoppelt werden.
Der f"ur das Koppeln n"otige Mechanismus wird weitestgehend 
vor dem Benutzer verborgen. Im folgenden Beispiel wird zuerst ein
Sechsring gezeichnet: \type{SIX,B}, die Position in Richtung der $x$-Achse
verschoben: \type{ADJ1} und ein F"unfring gezeichnet: \type{FIVE,ROT3,B}.

\startbuffer
\startchemie[kader=aan,breedte=4000,rechts=3000]
  \chemie[SIX,B,ADJ1,FIVE,ROT3,B]
\stopchemie
\stopbuffer

\startfiguurtekst
  [links][]
  {geen}
  {\haalbuffer}
\voorbeeld
\starttypen
\startchemie[rahmen=an,breite=4000,rechts=3000]
  \chemie[SIX,B,ADJ1,FIVE,ROT3,B]
\stopchemie
\stoptypen
\stopfiguurtekst

\startbuffer
\startchemie[kader=aan,breedte=5000,rechts=4500]
  \chemie[SIX,ROT2,B,R6,SUB1,FIVE,B,R4]
\stopchemie
\stopbuffer

\startfiguurtekst
  [links][]
  {geen}
  {\haalbuffer}
\voorbeeld
\starttypen
\startchemie[rahmen=an,breite=5000,rechts=4500]
  \chemie[SIX,ROT2,B,R6,SUB1,FIVE,B,R4]
\stopchemie
\stoptypen
\stopfiguurtekst

Eine Kopplung zu einer andersartigen Struktur geschieht mit \type{ADJ}. Wobei
die einzelnen Strukturen ggf. mit \type{ROT} gedreht werden m"ussen. Soll eine
Struktur nicht direkt, sondern "uber eine Bindung verkn"upft werden, verwendet
man \type{SUB}. Drehungen sind in 90 Grad Schritten im Uhrzeigersinn,
Verschiebungen durch \type{ADJ} und \type{SUB} in Richtung der Achsen m"oglich.

Die folgenden Beispiele illustrieren noch, da"s die Abmessungen der kleinen
Strukturen an die der gro"sen angepa"st sind. Die Struktur \type{EIGHT} hat 
"ubrigens etwas weniger M"oglichkeiten als zum Beispiel \type{SIX}.

\startbuffer
\startchemie[breedte=6000,links=1500,kader=aan]
  \chemie[EIGHT,B,MOV1,B]
\stopchemie
\stopbuffer

\startfiguurtekst
  [links][]
  {geen}
  {\haalbuffer}
\voorbeeld
\starttypen
\startchemie[breite=6000,links=1500,rahmen=an]
  \chemie[EIGHT,B,MOV1,B]
\stopchemie
\stoptypen
\stopfiguurtekst

\startbuffer
\startchemie[breedte=6000,links=1500,kader=aan]
  \chemie[EIGHT,B,ADJ1,SIX,B]
\stopchemie
\stopbuffer

\startfiguurtekst
  [links][]
  {geen}
  {\haalbuffer}
\voorbeeld
\starttypen
\startchemie[breite=6000,links=1500,rahmen=an]
  \chemie[EIGHT,B,ADJ1,SIX,B]
\stopchemie
\stoptypen
\stopfiguurtekst

\startbuffer
\startchemie[breedte=6000,links=1500,kader=aan]
  \chemie[EIGHT,B,ADJ1,FIVE,ROT3,B]
\stopchemie
\stopbuffer

\startfiguurtekst
  [links][]
  {geen}
  {\haalbuffer}
\voorbeeld
\starttypen
\startchemie[breite=6000,links=1500,rahmen=an]
  \chemie[EIGHT,B,ADJ1,FIVE,ROT3,B]
\stopchemie
\stoptypen
\stopfiguurtekst


Inzwischen haben Sie sicher festgestellt, da"s die Reihenfolge der Parameter
wichtig ist. Die empfohlene Reihenfolge lautet wie folgt: 

\starttypen
\chemie
  [Struktur,                             % SIX, FIVE, ...
   Bindungen innerhalb der Struktur,     % B, C, EB, ...
   Bindungen au�erhalb der Struktur,     % R, DR, ...
   Platz der Atome,                      % Z
   Platz der Substituenten]              % RZ, -RZ, ...
  [Atome,
   Substituenten]
\stoptypen

Das Verbinden von Strukturen kommt in der Regel nach dem Drehen. Obwohl
dies vielleicht nicht sofort zu verstehen ist, hat dies durchaus Systematik.
Der Proze"s h"atte zu dem vereinfacht werden k"onnen. 
Wir experimentierten mit einem mehr oder weniger automatischen
Drehmechanismus, aber solche versteckten Drehungen k"onnen
zu unerwarteten Ergebnissen f"uhren, so da"s wir davon Abstand nahmen.
Zudem ist es einfacher eine Struktur zu erstellen und sie dann erst zu drehen.
So  sollte eine verkn"upfte Struktur besser erst als Unterstruktur 
definiert werden, und dann als Ganzes gedreht werden.

Oft ist ein Sechsring  mit einem oder mehreren \type{ONE}s verbunden.
Solche Kohlenstoffketten werden oft im Zigzack dargestellt.
In einem solchen Fall benutzen wir \type{DIR}. 

\startbuffer
\startchemie
  [schaal=klein,breedte=6000,hoogte=6000,kader=aan]
  \chemie[SIX,SB2356,DB14,Z2346,SR36,RZ36]
    [C,N,C,C,H,H_2]
  \chemie[PB:Z1,ONE,Z0,DIR8,Z0,SB24,DB7,Z27,PE]
    [C,C,CH_3,O]
  \chemie[PB:Z5,ONE,Z0,DIR6,Z0,SB24,DB7,Z47,PE]
    [C,C,H_3C,O]
  \chemie[SR24,RZ24]
    [CH_3,H_3C]
\stopchemie
\stopbuffer

\startfiguurtekst
  [links][]
  {geen}
  {\haalbuffer}
\voorbeeld
\starttypen
\startchemie
  [format=klein,breite=6000,hoehe=6000,rahmen=an]
  \chemie[SIX,SB2356,DB14,Z2346,SR36,RZ36]
    [C,N,C,C,H,H_2]
  \chemie[PB:Z1,ONE,Z0,DIR8,Z0,SB24,DB7,Z27,PE]
    [C,C,CH_3,O]
  \chemie[PB:Z5,ONE,Z0,DIR6,Z0,SB24,DB7,Z47,PE]
    [C,C,H_3C,O]
  \chemie[SR24,RZ24]
    [CH_3,H_3C]
\stopchemie
\stoptypen
\stopfiguurtekst

Damit die Bindungen der Kette die gleiche L"ange haben, kann man jene
als Substruktur positionieren. Dazu verwenden wir die Befehle \type{PB}
und \type{PE}.

\plaatstabel
{Positioneren.}
\starttabel[ | r T | l w(4cm) | p(6cm) | ]
\HL
\VL PB:.. \VL Picture Begin \VL Beginn einer Unterstruktur\VL\FR
\VL PE    \VL Picture End   \VL Ende einer Unterstruktur\VL\LR
\HL
\stoptabel

Direkt nach \type{PB} wird angegeben, wohin die Struktur
platziert wird. Es ist wichtig daran zu denken, da"s das unmittelbar
folgende Atom zentriert wird. Darum ist der Gebrauch dem eines Zentralatom
vorzuziehen.

Der Ansto"s zum Einf"uhren dieser Befehle lag darin begr"undet, da"s
eine Struktur auf eine optisch akzeptable Weise gezeichnet werden sollte.
Es gibt jedoch verschiedene M"oglich\-keiten solche Strukturen zu zeichen.
Die folgende Variante liefert auch das Resultat: 

\startbuffer
\startchemie
    [schaal=klein,breedte=6000,hoogte=6000,kader=aan]
  \chemie[SIX,SB2356,DB14,Z36,SR36,RZ36][N,C,H,H_2]
  \chemie[PB:Z1,ONE,Z0,DIR8,Z0,SB24,DB7,Z27,PE][C,C,CH_3,O]
  \chemie[PB:Z5,ONE,Z0,DIR6,Z0,SB24,DB7,Z47,PE][C,C,H_3C,O]
  \chemie[PB:Z2,ONE,Z0,DIR2,SB6,CZ0,PE][C,CH_3]
  \chemie[PB:Z4,ONE,Z0,DIR4,SB8,CZ0,PE][C,H_3C]
\stopchemie
\stopbuffer

\startfiguurtekst
  [links][]
  {geen}
  {\haalbuffer}
\voorbeeld
\starttypen
\startchemie
    [format=klein,breite=6000,hoehe=6000,rahmen=an]
  \chemie
    [SIX,SB2356,DB14,Z36,SR36,RZ36][N,C,H,H_2]
  \chemie
    [PB:Z1,ONE,Z0,DIR8,Z0,SB24,DB7,Z27,PE][C,C,CH_3,O]
  \chemie
    [PB:Z5,ONE,Z0,DIR6,Z0,SB24,DB7,Z47,PE][C,C,H_3C,O]
  \chemie
    [PB:Z2,ONE,Z0,DIR2,SB6,CZ0,PE][C,CH_3]
  \chemie
    [PB:Z4,ONE,Z0,DIR4,SB8,CZ0,PE][C,CH_3]
\stopchemie
\stoptypen
\stopfiguurtekst

Am effizientesten wird diese Struktur aber wie folgt definiert. Die letzte
L"osung ist aber bez"uglich der Typographie nicht unbedingt die Beste.

\startbuffer
\startchemie
    [schaal=klein,breedte=6000,hoogte=6000,kader=aan]
  \chemie[SIX,SB2356,DB14,Z,SR36,RZ36,SR1245,RZ24]
    [C,C,N,C,C,C,H,H_2,CH_3,H_3C]
  \chemie[PB:RZ1,ONE,Z0,SB2,DB7,Z27,PE]
    [C,CH_3,O]
  \chemie[PB:RZ5,ONE,Z0,SB4,DB7,Z47,PE]
    [C,H_3C,O]
\stopchemie
\stopbuffer

\startfiguurtekst
  [links][]
  {geen}
  {\haalbuffer}
\voorbeeld
\starttypen
\startchemie
    [format=klein,breite=6000,hoehe=6000,rahmen=an]
  \chemie[SIX,SB2356,DB14,Z,SR36,RZ36,SR1245,RZ24]
    [C,C,N,C,C,C,H,H_2,CH_3,H_3C]
  \chemie[PB:RZ1,ONE,Z0,SB2,DB7,Z27,PE]
    [C,CH_3,O]
  \chemie[PB:RZ5,ONE,Z0,SB4,DB7,Z47,PE]
    [C,H_3C,O]
\stopchemie
\stoptypen
\stopfiguurtekst

Es ist vielleicht allen aufgefallen, da"s die Gr"o"se der Struktur durch
die Substituenten festgestellt wird. Dabei z"ahlen Ketten, abgesehen von
dem ersten Atom, dazu. Dies f"ordert einen konstanten Aufbau.

Die Unterschiede zwischen \type{SUB} und \type{PB} sind gering, sollten
aber nicht vernachl"assigt werden. Vergleichen Sie beispielsweise folgende
Formeln:

\startbuffer
\startchemie
  [breedte=passend,kader=aan,schaal=klein]
  \chemie
    [SIX,ROT2,B,C,R236,RZ23,
     SUB1,ONE,OFF1,Z0,4OFF1,SB1,Z1]
    [HO,HO,CHOH,CH_2NH_2]
\stopchemie
\stopbuffer

\startfiguurtekst
  [links][]
  {geen}
  {\haalbuffer}
\voorbeeld
\starttypen
\startchemie
  [breite=passend,rahmen=an,format=klein]
  \chemie
    [SIX,ROT2,B,C,R236,RZ23,
     SUB1,ONE,OFF1,Z0,4OFF1,SB1,Z1]
    [HO,HO,CHOH,CH_2NH_2]
\stopchemie
\stoptypen
\stopfiguurtekst

\startbuffer
\startchemie
  [breedte=passend,kader=aan,schaal=klein]
  \chemie
    [SIX,ROT2,B,C,R236,RZ23,
     PB:RZ6,ONE,Z0,3OFF1,SB1,Z1,PE]
    [HO,HO,CHOH,CH_2NH_2]
\stopchemie
\stopbuffer

\startfiguurtekst
  [links][]
  {geen}
  {\haalbuffer}
\voorbeeld
\starttypen
\startchemie
  [breite=passend,rahmen=an,format=klein]
  \chemie
    [SIX,ROT2,B,C,R236,RZ23,
     PB:RZ6,ONE,Z0,3OFF1,SB1,Z1,PE]
    [HO,HO,CHOH,CH_2NH_2]
\stopchemie
\stoptypen
\stopfiguurtekst

Das Verwenden von \type{PB:} ist etwas umst"andlicher, f"uhrt aber zu einem
besseren Ergebnis. Wie man "ubrigens hier deutlich sieht, mu"s man in diesem
Fall selbst die Breite der Struktur angeben. Dies kommt daher, da"s die
Abmessungen der so angegebenden Substituenten nicht ber"ucksichtigt werden k"onnen.
Ansonsten w"urden die Strukturen inkonsistente Abmessungen bekommen und
das Kombinieren erschweren.

Gehen wir zuerst auf die Bedeutung von \type{OFF} ein. Bei \type{ONE}
kann \type{Z0} aus mehr als einem Atom bestehen. In diesem Fall ist
der reservierte Platz zu klein. Wenn vor \type{Z0} mehr Platz n"otig
ist, dann k"onnen die Bindungen~1, 2 und~8 mit dem Parameter \type{OFF}
(\citaat{Offset}) verschoben werden. Wie wir den Parameter verwenden,
ist dem untenstehenden Beispiel zu entnehmen.

\startbuffer
\startchemie
  [breedte=passend,formaat=klein,schaal=klein,kader=aan]
  \chemie
    [SIX,B,C,ADJ1,
     FIVE,ROT3,SB34,+SB2,-SB5,Z345,DR35,SR4,CRZ35,SUB1,
     ONE,OFF1,SB258,Z0,Z28]
    [C,N,C,O,O,
     CH,COOC_2H_5,COOC_2H_5]
\stopchemie
\stopbuffer

\startfiguurtekst
  [links][]
  {geen}
  {\haalbuffer}
\voorbeeld
\starttypen
\startchemie
  [breite=passend,groesse=klein,format=klein,rahmen=an]
  \chemie
    [SIX,B,C,ADJ1,
     FIVE,ROT3,SB34,+SB2,-SB5,Z345,DR35,SR4,CRZ35,SUB1,
     ONE,OFF1,SB258,Z0,Z28]
    [C,N,C,O,O,
     CH,COOC_2H_5,COOC_2H_5]
\stopchemie
\stoptypen
\stopfiguurtekst

Ein Verschieben schafft Platz f"ur ein weiteres Zeichen. Wenn wir mehr Platz
ben"otigen, k"onnen wir beispielsweise \type{3OFF1} verwenden.
Solche auf den ersten Blick recht komplizierten Definitionen sollten
am Besten schrittweise aufgebaut werden und dann als Ganzes eingef"ugt
werden.

Es taucht hier auch noch ein neuer Parameter auf: \type{CRZ}. Dieser Parameter
kann dazu benutzt werden, um ein Atom oder Molek"ul mit der normalen Bindungsl"ange
zu platzieren, was hier gew"unscht ist. Dasselbe h"atten wir mit dem Parameter
\type{RZ} erreichen k"onnen, wenn wir mittels des zweiten Arguments den
Zwischenraum angepa"st h"att: \type{{\,O}} anstelle von \type{O}.

Zeigen wir noch ein weiteres Beispiel, zu dem wir noch einen Rat geben wollen:
Mit der Auswahl {\em eines} Verfahrens dient man der Konsistenz des 
bearbeiteten Dokuments.


\startbuffer
\startchemie
  [breedte=passend,hoogte=passend,kader=aan,
   schaal=klein]
  \chemie
    [ONE,SB15,DB7,Z057,3OFF1,MOV1,Z0,3OFF1,MOV1,
     Z017,SB1357,MOV3,Z0,MOV3,SB1357,Z013,3OFF5,
     MOV5,Z0,3OFF5,SB5,Z5]
    [C,H_2N,NH,(CH_2)_3,C,COOH,H,\SL{NH},C,COOH,H,
     (CH_2)_2,HOOC]
\stopchemie
\stopbuffer

\startfiguurtekst
  [links][]
  {geen}
  {\haalbuffer}
\voorbeeld
\starttypen
\startchemie
  [breite=passend,hoehe=passend,rahmen=an,
   format=klein]
  \chemie
    [ONE,SB15,DB7,Z057,3OFF1,MOV1,Z0,3OFF1,MOV1,
     Z017,SB1357,MOV3,Z0,MOV3,SB1357,Z013,3OFF5,
     MOV5,Z0,3OFF5,SB5,Z5]
    [C,H_2N,NH,(CH_2)_3,C,COOH,H,\SL{NH},C,COOH,H,
     (CH_2)_2,HOOC]
\stopchemie
\stoptypen
\stopfiguurtekst

\startbuffer
\startchemie
  [breedte=passend,hoogte=passend,kader=aan,
   schaal=klein]
  \chemie [ONE,SB15,DB7,Z057,3OFF1][C,H_2N,NH]
  \chemie [MOV1,Z0,3OFF1]          [(CH_2)_3]
  \chemie [MOV1,Z017,SB1357]       [C,COOH,H]
  \chemie [MOV3,Z0]                [\SL{NH}]
  \chemie [MOV3,SB1357,Z013,3OFF5] [C,COOH,H]
  \chemie [MOV5,Z0,3OFF5,SB5,Z5]   [(CH_2)_2,HOOC]
\stopchemie
\stopbuffer

\startfiguurtekst
  [links][]
  {geen}
  {\haalbuffer}
\voorbeeld
\starttypen
\startchemie
  [breite=passend,hoehe=passend,rahmen=an,
   format=klein]
  \chemie [ONE,SB15,DB7,Z057,3OFF1][C,H_2N,NH]
  \chemie [MOV1,Z0,3OFF1]          [(CH_2)_3]
  \chemie [MOV1,Z017,SB1357]       [C,COOH,H]
  \chemie [MOV3,Z0]                [\SL{NH}]
  \chemie [MOV3,SB1357,Z013,3OFF5] [C,COOH,H]
  \chemie [MOV5,Z0,3OFF5,SB5,Z5]   [(CH_2)_2,HOOC]
\stopchemie
\stoptypen
\stopfiguurtekst

\startbuffer
\startchemie
  [breedte=passend,hoogte=passend,kader=aan,
   schaal=klein]
  \chemie
    [ONE,Z0,SAVE,MOV7,SB1357,Z017,3OFF5,MOV5,Z0,
     3OFF5,MOV5,SB15,DB7,Z057,RESTORE,
     MOV3,SB1357,Z013,MOV5,3OFF5,Z0,6OFF5,SB5,Z5]
    [\SL{NH},C,COOH,H,(CH_2)_3,C,H_2N,NH,C,COOH,H,
     (CH_2)_2,HOOC]
\stopchemie
\stopbuffer

\startfiguurtekst
  [links][]
  {geen}
  {\haalbuffer}
\voorbeeld
\starttypen
\startchemie
  [breite=passend,hoehe=passend,rahmen=an,
   format=klein]
  \chemie
    [ONE,Z0,SAVE,MOV7,SB1357,Z017,3OFF5,MOV5,Z0,
     3OFF5,MOV5,SB15,DB7,Z057,RESTORE,
     MOV3,SB1357,Z013,MOV5,3OFF5,Z0,6OFF5,SB5,Z5]
    [\SL{NH},C,COOH,H,(CH_2)_3,C,H_2N,NH,C,COOH,H,
     (CH_2)_2,HOOC]
\stopchemie
\stoptypen
\stopfiguurtekst

\startbuffer
\startchemie
  [breedte=passend,hoogte=passend,kader=aan,
   schaal=klein]
  \chemie
    [ONE,Z0,MOV7,SB1357,Z017,3OFF5,MOV5,Z0,
     3OFF5,MOV5,SB15,DB7,Z057,MOV0,MOV3,SB1357,
     Z013,MOV5,3OFF5,Z0,6OFF5,SB5,Z5]
    [\SL{NH},C,COOH,H,(CH_2)_3,C,H_2N,NH,C,COOH,H,
     (CH_2)_2,HOOC]
\stopchemie
\stopbuffer

\startfiguurtekst
  [links][]
  {geen}
  {\haalbuffer}
\voorbeeld
\starttypen
\startchemie
  [breite=passend,hoehe=passend,rahmen=an,
   format=klein]
  \chemie
    [ONE,Z0,MOV7,SB1357,Z017,3OFF5,MOV5,Z0,
     3OFF5,MOV5,SB15,DB7,Z057,MOV0,MOV3,SB1357,
     Z013,MOV5,3OFF5,Z0,6OFF5,SB5,Z5]
    [\SL{NH},C,COOH,H,(CH_2)_3,C,H_2N,NH,C,COOH,H,
     (CH_2)_2,HOOC]
\stopchemie
\stoptypen
\stopfiguurtekst


Was vorzuziehen ist, h"angt von dem Geschmack der Benutzers und
nat"urlich von der gew"unschten typografischen Qualit"at ab. Kommen wir nun
zum Gebrauch von \type{SAVE} und \type{RESTORE}. Mit diesem Parameter
werden Zustandsvariablen gespeichert und wieder abgerufen.

Zum Schlu"s zeigen wir noch einmal ein Beispiel f"ur eine Kombination von
\type{SIX} mit \type{FIVE}. Es zeigt auch den Gebrauch von \type{SS}.

\startbuffer
\startchemie
    [breedte=passend,hoogte=passend,kader=aan]
  \chemie
    [SIX,DB135,SB246,Z,SR6,RZ6][C,C,N,\SR{HC},N,C,NH_2]
  \chemie
    [SIX,MOV1,DB1,SB23,SS6,Z1..5,SR3,RZ3][N,\SL{CH},N,C,C,H]
\stopchemie
\stopbuffer

\startfiguurtekst
  [links][]
  {geen}
  {\haalbuffer}
\voorbeeld
\starttypen
\startchemie
    [breite=passend,hoehe=passend,rahmen=an]
  \chemie
    [SIX,DB135,SB246,Z,SR6,RZ6][C,C,N,\SR{HC},N,C,NH_2]
  \chemie
    [SIX,MOV1,DB1,SB23,SS6,Z1..5,SR3,RZ3][N,\SL{CH},N,C,C,H]
\stopchemie
\stoptypen
\stopfiguurtekst

\hoofdstuk{Sonstiges Pla\-tzieren um Atome}

Wir k"onnen Strukturformeln mit Symbolen und Texten versehen, zum
Beispiel:

\startbuffer
\startchemie
  [hoogte=4500,boven=1250,breedte=passend,kader=aan]
  \ondertekst
    {2,2-Dimethyl-3-etylpentan}
  \chemie
    [ONE,Z3570,SB1357]
    [CH_3,\T{1}{H_3C},CH_3,\SR{\LT{2}{C}}]
  \chemie
    [MOV1,OFF1,Z0,SB3]
    [\T{3}{CH}]
  \chemie
    [MOV3,Z0,SB3,MOV3,Z0,MOV7,MOV7]
    [CH_2,CH_3]
  \chemie
    [OFF1,SB1,MOV1,OFF1,Z0,2OFF1,SB1,Z1]
    [\T{4}{CH_2},\T{5}{CH_3}]
\stopchemie
\stopbuffer

\startfiguurtekst
  [links][]
  {geen}
  {\haalbuffer}
\voorbeeld
\starttypen
\startchemie
  [hoehe=4500,oben=1250,breite=passend,rahmen=an]
  \textunter
    {2,2-Dimethyl-3-etylpentan}
  \chemie
    [ONE,Z3570,SB1357]
    [CH_3,\T{1}{H_3C},CH_3,\SR{\LT{2}{C}}]
  \chemie
    [MOV1,OFF1,Z0,SB3]
    [\T{3}{CH}]
  \chemie
    [MOV3,Z0,SB3,MOV3,Z0,MOV7,MOV7]
    [CH_2,CH_3]
  \chemie
    [OFF1,SB1,MOV1,OFF1,Z0,2OFF1,SB1,Z1]
    [\T{4}{CH_2},\T{5}{CH_3}]
\stopchemie
\stoptypen
\stopfiguurtekst

Neben \type{\T} gibt es noch eine Reihe anderer Parameter. Die Anzahl der
den Parametern "ubergebenen Argumenten ist optional.  So k"onnen die 
Oxidationszahlen (r"omische Ziffern) mit f"uhrenden \type{\+} und \type{\-}, oder
mit mit~\type{\1} bis~\type{\7} dargestellt werden.

\plaatstabel
{Sonstiges Platzieren um Atome: Oxidationszahlen}
\starttabel[ | l | l | ]
\HL
\VL \type{\+{nummer}} \VL positive r\"omische Oxidationszahl       \VL\FR
\VL \type{\-{nummer}} \VL negative r\"omische Oxidationszahl       \VL\MR
\VL \type{\1}         \VL die r\"omische Zahl 1 (ohne $+$ oder $-$)\VL\MR
\VL \type{\7}         \VL gleiches f\"ur 2, 3, 4, 5, 6 und 7       \VL\LR
\HL
\stoptabel

Die Oxidationszahlen werden "uber dem zugeh"origen Atom zentriert.
So ergibt: 

\startbuffer
\plaatsformule
  \startformule
    \chemie{S}+\chemie{O_2}
    \chemie{GIVES}
    \chemie{\+{4}{S}\-{2}{O_2}}
  \stopformule
\stopbuffer

\starttypen
$$
  \chemie{S}+\chemie{O_2}
  \chemie{GIVES}
  \chemie{\+{4}{S}\-{2}{O_2}}
$$
\stoptypen

dies:

\haalbuffer

Wir k"onnen schier endlose Polymerketten mit \type{\[} und~\type{\]} zeichnen.
Hierbei kann beiden Parametern optional Argumente "ubergeben werden, wie im
Beispiel f"ur Polytetrafluorethylen (PTFE, Handelsnamen sind Teflon und Hostaflon TF):

\startbuffer
\startchemie[breedte=passend,kader=aan]
  \chemie
    [ONE,ZT5,SB5,OFF1,Z0,,OFF1,SB1,MOV1,SB5,OFF1,Z0,OFF1,SB1,ZT1]
    [\[,CF_2,CF_2,\]{\sl n}]
\stopchemie
\stopbuffer

\startfiguurtekst
  [links][]
  {geen}
  {\haalbuffer}
\voorbeeld
\starttypen
\startchemie[breite=passend,rahmen=an]
  \chemie
    [ONE,ZT5,SB5,OFF1,Z0,,OFF1,SB1,MOV1,SB5,OFF1,Z0,OFF1,SB1,ZT1]
    [\[,C,C,C,\]{\sl n}]
\stopchemie
\stoptypen
\stopfiguurtekst

\plaatstabel
{Sonstiges platzieren um Atome: Polymerklammern.}
\starttabel[ | l | l | l | ]
\HL
\VL \type{\[{unten}} \VL \type{\[{unten}{oben}} \VL linke Polymerklammer (Fortf\"uhren nach rechts) \VL\FR
\VL \type{\]{unten}} \VL \type{\]{unten}{oben}} \VL rechte Polymerklammer (Fortf\"uhren nach links) \VL\LR
\HL
\stoptabel

Nat"urlich k"onnen wir Text links, rechts, oben oder unten platzieren.
Wir k"onnen vor \type{\L}, \type{\R}, \type{\T} und \type{\B}, um den 
Abstand zu verkleinen, den Parameter \type{\X} angeben.

\plaatstabel
{Sonstiges platzieren um Atome: Text.}
\starttabel[ | l | l | ]
\HL
\VL \type{\L{text}} \VL Text links  platzieren \VL\FR
\VL \type{\R{text}} \VL Text rechts platzieren \VL\MR
\VL \type{\T{text}} \VL Text oben   platzieren \VL\MR
\VL \type{\B{text}} \VL Text unten  platzieren \VL\LR
\HL
\stoptabel

Neben diesen vier Parametern k"onnen wir auch alle nun folgenden
Kombinationen verwenden, einschlie"slich der zentrierten Varianten.

\leavevmode\hbox to \hsize
  {\def\sample#1{\chemie{#1{\oplus}{\ruledhbox{\ttx\string#1}}}}%
   \sample\TL\hss\sample\L \hss\sample\LC\hss\sample\BL\hss
   \sample\TR\hss\sample\R \hss\sample\RC\hss\sample\BR\hss
   \sample\LT\hss\sample\T \hss\sample\RT\hss\sample\LB\hss
   \sample\B \hss\sample\RB}

\leavevmode\hbox to \hsize
  {\def\sample#1{\chemie{\X#1{\oplus}{\ruledhbox{\ttx\string\X\string#1}}}}%
   \sample\TL\hss\sample\L \hss\sample\LC\hss\sample\BL\hss
   \sample\TR\hss\sample\R \hss\sample\RC\hss\sample\BR\hss
   \sample\LT\hss\sample\T \hss\sample\RT\hss\sample\LB\hss
   \sample\B \hss\sample\RB}

Eine andere M"oglichkeit ist das Ausrichten von Text, was in der 
\TEX||Terminologie {\em smashed} Text genannt wird.

\plaatstabel
{Sonstiges platzieren um Atome: Ausrichtung der Atome}
\starttabel[ | l | l | ]
\HL
\VL \type{\SL{text}} \VL linksb\"undig  \VL\FR
\VL \type{\SM{text}} \VL zentriert      \VL\MR
\VL \type{\SR{text}} \VL rechtsb\"undig \VL\LR
\HL
\stoptabel

Die erste Variante wird nun als Beispiel verwendet, die anderen lassen sich
sinngem"a"s ableiten. Der Text wird linksb"undig ausgegeben:

\startbuffer
\startchemie[kader=aan]
  \chemie[SIX,B,R36,RZ36][\SL{COOH},\SL{COOH}]
\stopchemie
\stopbuffer

\startfiguurtekst
  [links][]
  {geen}
  {\haalbuffer}
\voorbeeld
\starttypen
\startchemie[rahmen=an]
  \chemie[SIX,B,R36,RZ36][\SL{COOH},\SL{COOH}]
\stopchemie
\stoptypen
\stopfiguurtekst

Wir k"onnen Text mit den Befehlen \type{\textueber}, \type{\textmitte} 
und \type{\textunter} "uber, unter oder in der Mitte der Strukturen  
platzieren. Die exakte Position des Textes h"angt von den Ausma"sen
der Strukturformel ab.

\startbuffer
\plaatsformule
  \startformule
    \startchemie[breedte=passend]
      \chemie[SIX,B,Z1,MOV1,B][\hbox{$\bullet$}]
      \boventekst{{\sl trans}-Decalin}
    \stopchemie
    \hskip 24pt
    \startchemie[breedte=passend]
      \chemie[SIX,B,Z12,MOV1,B][\hbox{$\bullet$},\hbox{$\bullet$}]
      \ondertekst{{\sl cis}-Decalin}
    \stopchemie
  \stopformule
\stopbuffer

\starttypen
$$
    \startchemie[breite=passend]
      \chemie[SIX,B,Z1,MOV1,B][\hbox{$\bullet$}]
      \textueber{{\sl trans}-Decalin}
    \stopchemie
    \hskip 24pt
    \startchemie[breite=passend]
      \chemie[SIX,B,Z12,MOV1,B][\hbox{$\bullet$},\hbox{$\bullet$}]
      \textunter{{\sl cis}-Decalin}
    \stopchemie
$$
\stoptypen

Beide {\em Decalin}||Formeln sehen wie folgt aus: 

\haalbuffer 

\hoofdstuk{Achsen}

Strukturen werden in einem abgegrenzten Bereich gezeichnet, der
aus Bequemlichkeit mit Achsen dargestellt werden kann. Die Abmessungen
der Achsen und deren Ursprung kann eingestellt werden. Zus"atzlich k"onnen
die Achsen angezeigt werden (z.\,B. zum Positionieren von Text) und
ein Rahmen kann gezeichent werden.

\startbuffer
\startchemie
  [assenstelsel=aan,
   breedte=6000,hoogte=4000]
\stopchemie
\stopbuffer

\startfiguurtekst
  [links][]{geen}
  {\haalbuffer}
\voorbeeld[vb:assenstelsel]
\starttypen
\startchemie
  [achsen=an,
   breite=6000,hoehe=4000]
\stopchemie
\stoptypen
\stopfiguurtekst

\startbuffer
\startchemie
  [assenstelsel=aan,
   links=2000,rechts=4000]
\stopchemie
\stopbuffer

\startfiguurtekst
  [links][]{geen}
  {\haalbuffer}
\voorbeeld
\starttypen
\startchemie
  [achsen=an,
   links=2000,rechts=4000]
\stopchemie
\stoptypen
\stopfiguurtekst

\startbuffer
\startchemie
  [assenstelsel=aan,
   breedte=6000,rechts=1000]
\stopchemie
\stopbuffer

\startfiguurtekst
  [links][]{geen}
  {\haalbuffer}
\voorbeeld
\starttypen
\startchemie
  [achsen=an,
   breite=6000,rechts=1000]
\stopchemie
\stoptypen
\stopfiguurtekst

\startbuffer
\startchemie
  [assenstelsel=aan,
   breedte=6000,boven=1000,onder=3000]
\stopchemie
\stopbuffer

\startfiguurtekst
  [links][]{geen}
  {\haalbuffer}
\voorbeeld
\starttypen
\startchemie
  [achsen=an,
   breite=6000,oben=1000,unten=3000]
\stopchemie
\stoptypen
\stopfiguurtekst

Die Abmessungen der gesammten Struktur werden den Abmessungen der
Achsen entnommen. Wenn jedoch \type{breite=passend} und/oder 
\type{hoehe=passend} angegeben wird, werden die Ausmessungen
entsprechend der wirklichen Gr"o"se der Struktur gesetzt.
Die Struktur wird, gem"a"s der Einstellungen, in den Text platziert:
nebeneinander, "ubereinander usw. \in{Beispiel}[vb:assenstelsel] zeigt
die Voreinstellung.

Innerhalb eines \type{\start}||\type{\stop}||Paars k"onnen \PICTEX||Makros
verwendet werden. Einige Vorsicht ist dabei nat"urlich geboten.

\hoofdstuk{Einstellungen}

Hinter \type{\startchemie} und \type{\stellechemieein} k"onnen verschiedene
Einstellungen angegeben werden.

\plaatstabel
{Einstellungen bei Strukturformeln.}
\starttabel[|lT|Tl|lT|]
\HL
\VL \bf Variable  \VL \bf Einstellungen          \VL \bf Vorgabe \VL\SR
\HL
\VL breite        \VL Zahl                       \VL 4000        \VL\FR
\VL hoehe         \VL Zahl                       \VL 4000        \VL\MR
\VL links         \VL Zahl                       \VL             \VL\MR
\VL rechts        \VL Zahl                       \VL             \VL\MR
\VL oben          \VL Zahl                       \VL             \VL\MR
\VL unten         \VL Zahl                       \VL             \VL\LR
\HL
\VL aufloesung    \VL Zahl                       \VL \type{\outputresolution} \VL\FR
\VL fliesstext    \VL 8pt 9pt 10pt usw.          \VL \type{\bodyfontsize}     \VL\MR
\VL schrift       \VL \type{\rm} \type{\bf} usw. \VL \type{\rm}               \VL\LR
\HL
\VL format         \VL Zahl klein mittel gross    \VL mittel     \VL\FR
\VL groesse       \VL klein mittel gross         \VL mittel      \VL\LR
\HL
\VL status        \VL start stop                 \VL start       \VL\FR
\VL option        \VL test                       \VL             \VL\LR
\HL
\VL achsen        \VL an aus                     \VL aus         \VL\FR
\VL rahmen        \VL an aus                     \VL aus         \VL\LR
\HL
\VL alternative   \VL 1 2                        \VL 1           \VL\SR
\HL
\VL offset        \VL HIGH LOW                   \VL LOW         \VL\SR
\HL
\stoptabel

Standardgem"a"s ist der Bereich der Achsen (H"ohe wie Breite) auf $-2000$ 
bis $+2000$ eingestellt. Der Ursprungspunkt \type{Z0} liegt dabei bei (0,0).
Andere Werte k"onnen mit Hilfe von \type{links}, \type{rechts}, 
\type{oben} und/oder \type{unten} in Kombination mit \type{breite}
und \type{hoehe} eingestellt werden.

Mit \type{groesse} kann die Abmessung der Buchstaben eingestellt werden.
Dabei werden die \TEX||Primitive \type{\textsize}, \type{\scriptsize} und
\type{\scriptscriptsize} verwendet. Man stellt mit \type{format} die
Abmessungen der Struktur selbst ein (\type{1..1000}). Die Gr"o"se wird durch
den \type{fliesstext} bestimmt. Die Schl"usselw"orter \type{klein}, 
\type{mittel} und \type{gross} sind auf einander abgestimmt.

Die Einstellung \type{fliesstext} wird zum Berechnen ben"utzt und hat 
auf den laufenden Text keinen Einflu"s. Innerhalb von \CONTEXT\ und
\LATEX\ ist er automatisch mit der Gr"o"se des Flie"stextes gekoppelt.

Sowohl in \TEX\ als auch in \CONTEXT\ k"onnen im mathematischen Modus
Befehle wie \type{\rm}, \type{\bf} und \type{\sl} verwendet werden.
Standardgem"a"s benutzt \PPCHTEX\ \type{\rm}. Mit der Einstellung 
\type{schrift} kann eine andere Einstellung gesetzt werden. In 
\in{Beispiel}[vb:letter] wurden die Substituenten {\sl geneigt} gesetzt.

\startbuffer
\startchemie[kader=aan,letter=\sl]
  \chemie[SIX,B,R,RZ][A,B,C,D,E,F]
\stopchemie
\stopbuffer

\startfiguurtekst
  [links][]
  {geen}
  {\haalbuffer}
\voorbeeld[vb:letter]
\starttypen
\startchemie[rahmen=an,schrift=\sl]
  \chemie[SIX,B,R,RZ][A,B,C,D,E,F]
\stopchemie
\stoptypen
\stopfiguurtekst

Die Einstellung der \type{schrift} gilt (vorl"aufig) sowohl f"ur Textformeln
als auch f"ur Strukturformeln. Die Indizies werden mit ver"andert, beispielsweise:
{\stelchemiein[letter=\rm]\chemie{CH_4}},{\stelchemiein[letter=\bf]\chemie{CH_4}} 
und {\stelchemiein[letter=\sl]\chemie{CH_4}}; dazu haben wir diese
Einstellungen verwendet: \type{\rm}, \type{\bf} und \type{\sl}.

Kursive ({\em italic,} \type{\it}) Formeln f"uhren zu gr"o"seren Zeilenabst"anden.
Innerhalb von \CONTEXT\ sind auch halbfettgeneigt ({\em bold||slanted,} \type{\bs})
und halbfettkursiv ({\em bold||italic\,} \type{\bi}) standardgem"a"s verf"ugbar.
Diese Befehle passen sich automatisch an den jeweiligen Stil an: 
{\stelchemiein[letter=\ss\sl]\chemie{CH_4}},{\stelchemiein[letter=\rm\sl]\chemie{CH_4}},
{\stelchemiein[letter=\tt\sl]\chemie{CH_4}} etc. (\type{\ss}, \type{\rm},
\type{\tt}). 

Nat"urlich k"onnen wir auch die Schriftart direkt einstellen, wie das
folgende Beispiel zeigt:

\startbuffer
\startchemie[kader=aan]
  \chemie[SIX,B,R,RZ][\tf a,\bf b,\it c,\sl d,\bi e,\bs f]
\stopchemie
\stopbuffer

\startfiguurtekst
  [links][]
  {geen}
  {\haalbuffer}
\voorbeeld
\starttypen
\startchemie[rahmen=an]
  \chemie[SIX,B,R,RZ][\tf a,\bf b,\it c,\sl d,\bi e,\bs f]
\stopchemie
\stoptypen
\stopfiguurtekst

Mit \type{status} kann die lange Rechenzeit -- auf Kosten der Qualit"at
-- verk"urzt werden. Die Variablen \type{rahmen} und \type{achsen} 
sind selbst erkl"arend. Mit der Option \type{test} kann ein Rahmen um 
den Text der Formel gezeichnet werden.\voetnoot{Innerhalb von \CONTEXT\ wird
dies zur Wiedergabe im visuellen Debugger verwendet. Dies zeigt die
Grundlinie als punktierte Linie an. Das Modul \type{supp-vis} wird 
"ubrigens auch in anderen Systemen verwendet.} Mit \type{alternative} 
stellen wir die Liniendicke ein. Standardgem"a"s ben"utzt \PPCHTEX\ 
einen 5~Punkte gro"sen Punkt. Wenn \type{alternative} 2 benutzt wird, wird
ein kleinerer Punkt verwendet, d.\,h. es werden d"unnere Linien gezeichnet.

\startbuffer
\startchemie[kader=aan,optie=test,variant=2]
  \chemie[SIX,B,R,RZ][R_1,R_2,R_3,R_4,R_5,R_6]
\stopchemie
\stopbuffer

\startfiguurtekst
  [links][]
  {geen}
  {\haalbuffer}
\voorbeeld[vb:test]
\starttypen
\startchemie[rahmen=an,option=test,alternative=2]
  \chemie[SIX,B,R,RZ][R_1,R_2,R_3,R_4,R_5,R_6]
\stopchemie
\stoptypen
\stopfiguurtekst

Das \type{offset} bezieht sich auf die Position von hoch|| und tiefgestelltem
Text. Bei \type{HIGH} werden die Indizies (tiefgestellter Text) etwas 
h"oher(\chemie{HIGH,H_2O}), bei \type{LOW} etwas tiefer platziert 
(\chemie{LOW,H_2O}).

Eine Struktur kann in verschieden Gr"o"sen gezeichnet werden. Die 
Einstellung geschieht mit \type{groesse} und \type{format}.

\startbuffer
\startchemie[kader=aan,schaal=klein,formaat=klein]
  \chemie[SIX,B,R,RZ][1,2,3,4,5,6]
\stopchemie
\stopbuffer

\startfiguurtekst
  [links][]
  {geen}
  {\haalbuffer}
\voorbeeld
\starttypen
\startchemie[rahmen=an,format=klein,groesse=klein]
  \chemie[SIX,B,R,RZ][1,2,3,4,5,6]
\stopchemie
\stoptypen
\stopfiguurtekst

\startbuffer
\startchemie[kader=aan,schaal=middel,formaat=middel]
  \chemie[SIX,B,R,RZ][1,2,3,4,5,6]
\stopchemie
\stopbuffer

\startfiguurtekst
  [links][]
  {geen}
  {\haalbuffer}
\voorbeeld
\starttypen
\startchemie[rahmen=an,format=mittel,groesse=mittel]
  \chemie[SIX,B,R,RZ][1,2,3,4,5,6]
\stopchemie
\stoptypen
\stopfiguurtekst

\startbuffer
\startchemie[kader=aan,schaal=groot,formaat=groot]
  \chemie[SIX,B,R,RZ][1,2,3,4,5,6]
\stopchemie
\stopbuffer

\startfiguurtekst
  [links][]
  {geen}
  {\haalbuffer}
\voorbeeld
\starttypen
\startchemie[rahmen=an,format=gross,groesse=gross]
  \chemie[SIX,B,R,RZ][1,2,3,4,5,6]
\stopchemie
\stoptypen
\stopfiguurtekst

Die \type{format} kann auch mit einer Zahl zwischen~1 und 1000 angegeben werden.
Die Werte der Schl"usselw"orter \type{klein}, \type{mittel} und \type{gross}
sind auf einander abgestimmt.

\hoofdstuk{Symbole}

F"ur das Setzen von Reaktionsgleichungen sind einige Symbole verf"ugbar.
Im folgenden Beispiel wird eine Gleichung dargestellt, die wie folgt 
definiert ist:

\startbuffer
\stelchemiein
  [formaat=klein,
   schaal=klein,
   breedte=passend,
   hoogte=5500,
   onder=1500]

\hbox
  {\startchemie
     \chemie[SIX,B,ER6,RZ6][O]
   \stopchemie
   \startchemie
     \chemie[SPACE,PLUS,SPACE]
   \stopchemie
   \startchemie
     \chemie[FIVE,ROT4,B125,+SB3,-SB4,Z4,SR4,RZ4][N,H]
   \stopchemie
   \startchemie
     \chemie[SPACE,GIVES,SPACE][?]
   \stopchemie
   \startchemie
     \chemie[SIX,B,EB6,R6,SUB4,FIVE,ROT4,B125,+SB3,-SB4,Z4][N]
   \stopchemie
   \startchemie
     \chemie[SPACE,PLUS,SPACE,CHEM][H_2O]
   \stopchemie}
\stopbuffer

\starttypen
\stellechemieein
  [groesse=klein,
   format=klein,
   breite=passend,
   hoehe=5500,
   unten=1500]

\hbox
  {\startchemie
     \chemie[SIX,B,ER6,RZ6][O]
   \stopchemie
   \startchemie
     \chemie[SPACE,PLUS,SPACE]
   \stopchemie
   \startchemie
     \chemie[FIVE,ROT4,B125,+SB3,-SB4,Z4,SR4,RZ4][N,H]
   \stopchemie
   \startchemie
     \chemie[SPACE,GIVES,SPACE][?]
   \stopchemie
   \startchemie
     \chemie[SIX,B,EB6,R6,SUB4,FIVE,ROT4,B125,+SB3,-SB4,Z4][N]
   \stopchemie
   \startchemie
     \chemie[SPACE,PLUS,SPACE,CHEM][H_2O]
   \stopchemie}
\stoptypen

Man kann \type{CHEM} dazu benutzen, Text (chemische Formeln)
zu setzen. Wenn man dagegen \type{TEXT} gebraucht, wird die
aktuelle Schrift verwendet.

Der Befehl \type{\hbox} ist n"otig, um die Struktur hintereinander zu setzen.
Die Symbole \type{GIVES} und \type{PLUS} sind selbsterkl"arend. Mit 
\type{SPACE} kann ein zus"atzlicher Zwischenraum eingef"uget werden.

\plaatsfiguur
  [hier]
  [fig:reactie]
  {geen}
  {\haalbuffer}

Ein Gleichgewicht kann mit \type{EQUILIBRIUM} wiedergegeben werden. "Uber
\type{EQUILIBRIUM} und \type{GIVES} kann ein Text gesetzt werden. Im 
Beispiel ist dies ein \citaat{?}.

Neben \type{GIVES} und \type{EQUILIBRIUM} ist auch \type{MESOMERIC} m"oglich.
Daneben k"onnen auch Komplexe mit Hilfe von 
\type{OPENCOMPLEX} und \type{CLOSECOMPLEX} gesetzt werden.

\hoofdstuk{Positionieren}

Beim Kombinieren von Molek"ulen, zum Beispiel mit \type{SUB}, werden einige
Gr"o"sen wie die Ausma"se, die Position, sowie der Mittelpunkt ver"andert.
Deshalb ist es m"oglich jene Gr"o"sen mit \type{SAVE} zu speichern und 
mit \type{RESTORE} wiederaufzurufen.

\plaatstabel
{Positionieren.}
\starttabel[ | r T | l w(4cm) | p(7cm) | ]
\HL
\VL SAVE    \VL Save Status    \VL Speichern der Positionsvariablen \VL\FR
\VL RESTORE \VL Restore Status \VL Wiederherstellen der Variablen   \VL\LR
\HL
\stoptabel

Wir verwenden die Parameter \type{SAVE} und \type{RESTORE} vor allem bei
Substituenten; bei Radikalen sollte man \type{PB} und \type{PE} benutzen,
die einen genaueren Positionierungsmechanismus verwenden.

\startbuffer
\definieerchemie[molecuul]
  {\chemie
     [ONE,Z0,SB1357,
      SAVE,SUB2,SIX,B,R6,C,RESTORE,
      MOV1,Z0,SB137,
      MOV1,Z0,SB37,
      MOV1]
     [C,C,C]}

\startchemie[breedte=passend,hoogte=passend]
  \chemie[molecuul,molecuul,molecuul]
\stopchemie
\stopbuffer

\starttypen
\definierechemie[molekuel]
  {\chemie
     [ONE,Z0,SB1357,
      SAVE,SUB2,SIX,B,R6,C,RESTORE,
      MOV1,Z0,SB137,
      MOV1,Z0,SB37,
      MOV1]
     [C,C,C]}

\startchemie[breite=passend,hoehe=passend]
  \chemie[molekuel,molekuel,molekuel]
\stopchemie
\stoptypen

Dieses Beispiel zeigt wie mehrere Molek"ule zu einem einzigen verbunden
werden k"onnen.

\plaatsfiguur
  [hier]
  [fig:save]
  {geen}
  {\haalbuffer}

Das n"achste Beispiel ist etwas fortgeschrittener und zeigt eine 
M"oglichkeit, Gleichungen zu setzten. Bei solchen Formeln mu"s man
die H"ohe der Formeln selbst einstellen.

\startbuffer
\plaatsformule
  \startformule
    \stelchemiein
      [breedte=passend,boven=2000,onder=2000,
       schaal=klein,formaat=klein]
    \startchemie
      \chemie
        [ONE,
         SAVE,
           Z0,SB7,SB3,SB1,MOV1,Z0,SB1,MOV1,Z0,DB8,CZ8,SB1,Z1,
         RESTORE,
         SAVE,
           SUB4,ONE,Z0,SB3,SB1,MOV1,Z0,SB1,MOV1,Z0,DB8,CZ8,SB1,Z1,
         RESTORE,
         SUB2,ONE,Z0,SB7,SB1,MOV1,Z0,SB1,MOV1,Z0,DB8,CZ8,SB1,Z1]
        [\SR{HC},O,C,O,C_{19}H_{39},
         \SR{H_{2}C},O,C,O,C_{17}H_{29},
         \SR{H_{2}C},O,C,O,C_{21}H_{41}]
    \stopchemie
    \startchemie
      \chemie[SPACE,PLUS,SPACE]
    \stopchemie
    \startchemie[rechts=600]
      \chemie[ONE,CZ0][3CH_{3}OH]
    \stopchemie
    \startchemie
      \chemie[SPACE,GIVES,SPACE,SPACE][H^+/H_2O]
    \stopchemie
    \startchemie
      \chemie
        [ONE,
         SAVE,
           Z0,SB7,SB3,SB1,Z1,
         RESTORE,
         SAVE,
           SUB4,ONE,Z0,SB3,SB1,Z1,
         RESTORE,
         SUB2,ONE,Z0,SB7,SB1,Z1]
        [\SR{HC},OH,
         \SR{H_{2}C},OH,
         \SR{H_{2}C},OH]
    \stopchemie
    \startchemie
      \chemie[SPACE,PLUS,SPACE]
    \stopchemie
    \startchemie
      \chemie
        [ONE,
         SAVE,
           Z0,DB8,CZ8,SB1,SB5,Z5,MOV1,Z0,SB1,Z1,
         RESTORE,
         SAVE,
           SUB4,ONE,Z0,DB8,CZ8,SB1,SB5,Z5,MOV1,Z0,SB1,Z1,
         RESTORE,
         SUB2,ONE,Z0,DB8,CZ8,SB1,SB5,Z5,MOV1,Z0,SB1,Z1]
        [C,O,C_{19}H_{39},O,CH_{3},
         C,O,C_{17}H_{29},O,CH_{3},
         C,O,C_{21}H_{41},O,CH_{3}]
    \stopchemie
  \stopformule
\stopbuffer
\starttypen
$$
    \stellechemieein
      [breite=passend,oben=2000,unten=2000,
       format=klein,groesse=klein]
    \startchemie
      \chemie
        [ONE,
         SAVE,
           Z0,SB7,SB3,SB1,MOV1,Z0,SB1,MOV1,Z0,DB8,CZ8,SB1,Z1,
         RESTORE,
         SAVE,
           SUB4,ONE,Z0,SB3,SB1,MOV1,Z0,SB1,MOV1,Z0,DB8,CZ8,SB1,Z1,
         RESTORE,
         SUB2,ONE,Z0,SB7,SB1,MOV1,Z0,SB1,MOV1,Z0,DB8,CZ8,SB1,Z1]
        [\SR{HC},O,C,O,C_{19}H_{39},
         \SR{H_{2}C},O,C,O,C_{17}H_{29},
         \SR{H_{2}C},O,C,O,C_{21}H_{41}]
    \stopchemie
    \startchemie
      \chemie[SPACE,PLUS,SPACE]
    \stopchemie
    \startchemie[rechts=600]
      \chemie[ONE,CZ0][3CH_{3}OH]
    \stopchemie
    \startchemie
      \chemie[SPACE,GIVES,SPACE,SPACE][H^+/H_2O]
    \stopchemie
    \startchemie
      \chemie
        [ONE,
         SAVE,
           Z0,SB7,SB3,SB1,Z1,
         RESTORE,
         SAVE,
           SUB4,ONE,Z0,SB3,SB1,Z1,
         RESTORE,
         SUB2,ONE,Z0,SB7,SB1,Z1]
        [\SR{HC},OH,
         \SR{H_{2}C},OH,
         \SR{H_{2}C},OH]
    \stopchemie
    \startchemie
      \chemie[SPACE,PLUS,SPACE]
    \stopchemie
    \startchemie
      \chemie
        [ONE,
         SAVE,
           Z0,DB8,CZ8,SB1,SB5,Z5,MOV1,Z0,SB1,Z1,
         RESTORE,
         SAVE,
           SUB4,ONE,Z0,DB8,CZ8,SB1,SB5,Z5,MOV1,Z0,SB1,Z1,
         RESTORE,
         SUB2,ONE,Z0,DB8,CZ8,SB1,SB5,Z5,MOV1,Z0,SB1,Z1]
        [C,O,C_{19}H_{39},O,CH_{3},
         C,O,C_{17}H_{29},O,CH_{3},
         C,O,C_{21}H_{41},O,CH_{3}]
    \stopchemie
$$
\stoptypen

Diese Definition h"atte man auch, auf Kosten der Lesbarkeit, k"urzer 
darstellen k"onnen (\type{SB731} statt \type{SB7,SB3,SB1}). Solche 
umfangreiche Strukturen sollte man am Besten in einer separaten Datei 
platzieren (eventuell in \PLAIN\ \TEX).

\haalbuffer

Nun geben wir noch zwei Beispiele bei denen ein Text unter der Struktur
platziert wird.

\startbuffer
\plaatsformule
  \startformule
    \stelchemiein
      [breedte=passend,boven=1500,onder=3500]
    \startchemie
      \chemie
        [ONE,Z0,DB1,SB3,SB7,Z7,MOV1,Z0,SB3,SB7,Z3,Z7,
         MOV0,SUB2,SIX,B,R6,C]
        [C,H,C,H,H]
      \ondertekst{Styreen}
    \stopchemie
    \quad\quad\quad
    \startchemie
      \chemie
        [ONE,Z0,DB1,SB3,SB7,Z3,Z7,
         MOV1,Z0,SB1,SB3,Z3,
         MOV1,Z0,DB1,SB3,Z3,
         MOV1,Z0,SB3,SB7,Z3,Z7]
        [C,H,H,C,H,C,H,C,H,H]
      \ondertekst{1,3-Butadieen}
   \stopchemie
  \stopformule
\stopbuffer

\starttypen
$$
    \stellechemieein
      [breite=passend,oben=1500,unten=3500]
    \startchemie
      \chemie
        [ONE,Z0,DB1,SB3,SB7,Z7,MOV1,Z0,SB3,SB7,Z3,Z7,
         MOV0,SUB2,SIX,B,R6,C]
        [C,H,C,H,H]
      \textunter{Styreen}
    \stopchemie
    \quad\quad\quad
    \startchemie
      \chemie
        [ONE,Z0,DB1,SB3,SB7,Z3,Z7,
         MOV1,Z0,SB1,SB3,Z3,
         MOV1,Z0,DB1,SB3,Z3,
         MOV1,Z0,SB3,SB7,Z3,Z7]
        [C,H,H,C,H,C,H,C,H,H]
      \textunter{1,3-Butadieen}
   \stopchemie
$$
\stoptypen

\haalbuffer

Der Gebrauch von \type{OFF} kommt etwas kurz. Die folgenden Beispiele lassen
aber erkennen, da"s exakte Verschiebungen m"oglich sind.

\startbuffer
\startchemie[hoogte=passend,kader=aan]
  \chemie[ONE,SB]
\stopchemie
\stopbuffer

\startfiguurtekst
  [links][]
  {geen}
  {\haalbuffer}
\voorbeeld
\starttypen
\startchemie[hoehe=passend,rahmen=an]
  \chemie[ONE,SB]
\stopchemie
\stoptypen
\stopfiguurtekst

\startbuffer
\startchemie[hoogte=passend,kader=aan]
  \chemie[ONE,3OFF1,SB]
\stopchemie
\stopbuffer

\startfiguurtekst
  [links][]
  {geen}
  {\haalbuffer}
\voorbeeld
\starttypen
\startchemie[hoehe=passend,rahmen=an]
  \chemie[ONE,3OFF1,SB]
\stopchemie
\stoptypen
\stopfiguurtekst

\startbuffer
\startchemie[hoogte=passend,kader=aan]
  \chemie[ONE,MOV1,SB]
\stopchemie
\stopbuffer

\startfiguurtekst
  [links][]
  {geen}
  {\haalbuffer}
\voorbeeld
\starttypen
\startchemie[hoehe=passend,rahmen=an]
  \chemie[ONE,MOV1,SB]
\stopchemie
\stoptypen
\stopfiguurtekst

\startbuffer
\startchemie[hoogte=passend,kader=aan]
  \chemie[ONE,3OFF1,MOV1,SB]
\stopchemie
\stopbuffer

\startfiguurtekst
  [links][]
  {geen}
  {\haalbuffer}
\voorbeeld
\starttypen
\startchemie[hoehe=passend,rahmen=an]
  \chemie[ONE,3OFF1,MOV1,SB]
\stopchemie
\stoptypen
\stopfiguurtekst

\startbuffer
\startchemie[hoogte=passend,kader=aan]
  \chemie[ONE,MOV1,3OFF1,SB]
\stopchemie
\stopbuffer

\startfiguurtekst
  [links][]
  {geen}
  {\haalbuffer}
\voorbeeld
\starttypen
\startchemie[hoehe=passend,rahmen=an]
  \chemie[ONE,MOV1,3OFF1,SB]
\stopchemie
\stoptypen
\stopfiguurtekst

\startbuffer
\startchemie[hoogte=passend,kader=aan]
  \chemie[ONE,MOV1,3OFF1,OFF0,SB]
\stopchemie
\stopbuffer

\startfiguurtekst
  [links][]
  {geen}
  {\haalbuffer}
\voorbeeld
\starttypen
\startchemie[hoehe=passend,rahmen=an]
  \chemie[ONE,MOV1,3OFF1,OFF0,SB]
\stopchemie
\stoptypen
\stopfiguurtekst

\startbuffer
\startchemie[hoogte=passend,kader=aan]
  \chemie[ONE,MOV1,3OFF1,MOV0,SB] 
\stopchemie
\stopbuffer

\startfiguurtekst
  [links][]
  {geen}
  {\haalbuffer}
\voorbeeld
\starttypen
\startchemie[hoehe=passend,rahmen=an]
  \chemie[ONE,MOV1,3OFF1,MOV0,SB] 
\stopchemie
\stoptypen
\stopfiguurtekst

\startbuffer
\startchemie[hoogte=passend,kader=aan]
  \chemie[ONE,MOV1,MOV0,SB]
\stopchemie
\stopbuffer

\startfiguurtekst
  [links][]
  {geen}
  {\haalbuffer}
\voorbeeld
\starttypen
\startchemie[hoehe=passend,rahmen=an]
  \chemie[ONE,MOV1,MOV0,SB]
\stopchemie
\stoptypen
\stopfiguurtekst

Das folgende Beispiel zeigt wie Komplexe dargestellt werden k"onnen.
Auch demonstriert es den Gebrauch von \type{RBT}. Normalerweise ist
der zus"atzliche Zwischenraum nicht n"otig, aber wir gebrauchen hier
-- der Befehl ist nicht sichtbar -- eine noch kleinere Schrift, damit
die Abbildung nicht zu breit wird.

\startbuffer
\startformule
\stelchemiein[schaal=klein,breedte=passend]
\startchemie
  \chemie[OPENCOMPLEX,SPACE]
\stopchemie
\startchemie
  \chemie[SIX,B,EB35,R6,-R1,+R1]
  \chemie[SIX,PB:RZ6,ONE,OFF1,Z0,EP57,PE][\SL{OH}]
  \chemie[SIX,-RZ1,+RZ1,RT2][H,Br,\oplus]
\stopchemie
\startchemie
  \chemie[SPACE,MESOMERIC,SPACE]
\stopchemie
\startchemie
  \chemie[SIX,B,EB25,R6,-R1,+R1]
  \chemie[SIX,PB:RZ6,ONE,OFF1,Z0,EP57,PE][\SL{OH}]
  \chemie[SIX,-RZ1,+RZ1,RT4][H,Br,\oplus]
\stopchemie
\startchemie
  \chemie[SPACE,MESOMERIC,SPACE]
\stopchemie
\startchemie
  \chemie[SIX,B,EB24,R6,-R1,+R1]
  \chemie[SIX,PB:RZ6,ONE,OFF1,Z0,EP57,PE][\SL{OH}]
  \chemie[SIX,-RZ1,+RZ1,RBT6][H,Br,\ \oplus] 
\stopchemie
\startchemie
  \chemie[SPACE,MESOMERIC,SPACE]
\stopchemie
\startchemie
  \chemie[SIX,B,EB24,ER6,-R1,+R1]
  \chemie[SIX,PB:RZ6,ONE,OFF1,Z0,EP5,ZT7,PE][\SL{OH},\oplus]
  \chemie[SIX,-RZ1,+RZ1][H,Br]
\stopchemie
\startchemie
  \chemie[SPACE,CLOSECOMPLEX]
\stopchemie
\stopformule
\stopbuffer

\starttypen
$$
\stelchemiein[format=klein,breite=passend]
\startchemie
  \chemie[OPENCOMPLEX,SPACE]
\stopchemie
\startchemie
  \chemie[SIX,B,EB35,R6,-R1,+R1]
  \chemie[SIX,PB:RZ6,ONE,OFF1,Z0,EP57,PE][\SL{OH}]
  \chemie[SIX,-RZ1,+RZ1,RT2][H,Br,\oplus]
\stopchemie
\startchemie
  \chemie[SPACE,MESOMERIC,SPACE]
\stopchemie
\startchemie
  \chemie[SIX,B,EB25,R6,-R1,+R1]
  \chemie[SIX,PB:RZ6,ONE,OFF1,Z0,EP57,PE][\SL{OH}]
  \chemie[SIX,-RZ1,+RZ1,RT4][H,Br,\oplus]
\stopchemie
\startchemie
  \chemie[SPACE,MESOMERIC,SPACE]
\stopchemie
\startchemie
  \chemie[SIX,B,EB24,R6,-R1,+R1]
  \chemie[SIX,PB:RZ6,ONE,OFF1,Z0,EP57,PE][\SL{OH}]
  \chemie[SIX,-RZ1,+RZ1,RBT6][H,Br,\ \oplus] 
\stopchemie
\startchemie
  \chemie[SPACE,MESOMERIC,SPACE]
\stopchemie
\startchemie
  \chemie[SIX,B,EB24,ER6,-R1,+R1]
  \chemie[SIX,PB:RZ6,ONE,OFF1,Z0,EP5,ZT7,PE][\SL{OH},\oplus]
  \chemie[SIX,-RZ1,+RZ1][H,Br]
\stopchemie
\startchemie
  \chemie[SPACE,CLOSECOMPLEX]
\stopchemie
$$
\stoptypen

Ohne den Gebrauch von \type{SPACE} w"urden die verschieden Elemente der
Abbildung zu dicht aneinander sein und sich ber"uhren. Das Optimieren 
der Abbildungen und "ahnlicher Reaktionsgleichungen ist oft ein
schrittweiser Proze"s.

\start
\switchnaarkorps[klein]
\haalbuffer
\stop

\hoofdstuk{Textformeln}

Neben dem Setzten von Strukturformeln wird auch das Setzen von 
Reaktionsgleichungen unterst"utzt. Der bereits beschriebene Befehl
\type{\chemie} hat noch drei weitere Varianten:

\starttypen
\chemie{Formel}
\chemie{Formel}{Text unten}
\chemie{Formel}{Text oben }{Text unten}
\stoptypen

Der Befehl pa"st sich an der Stelle im Dokument an; d.\,h. es 
werden die folgenden Modi unterschiedlich behandelt:

\startopsomming[opelkaar]
\som Textmodus
\som mathematischer Textmodus
\som mathematischer Displaymodus
\stopopsomming

Wenn dieser Befehl im laufenden Text gegeben wird, wird automatisch
\type{$ $} um den Befehl plaziert. So liefert \type{\chemie{NH_4^+}} die
Formel \chemie{NH_4^+}.

Das selbe Ergebnis erh"alt man, wenn man den Befehl zwischen \type{$ $}
plaziert. In beiden F"allen wird das zweite und dritte Argument weggelassen.
Wenn wir den Befehl zwischen \type{$$ $$} (oder in \CONTEXT\ \type{\startformel}
und \type{\stopformel}) platzieren, dann sind die zus"atzlichen Argumente
funktionsf"ahig.

\startbuffer
\plaatsformule
  \startformule
    \chemie{2H_2} \chemie{PLUS} \chemie{O_2}
    \chemie{GIVES} \chemie{2H_2O}
  \stopformule
\stopbuffer

\starttypen
$$
 \chemie{2H_2} \chemie{PLUS} \chemie{O_2}
 \chemie{GIVES} \chemie{2H_2O}
$$
\stoptypen

liefert:

\haalbuffer

Die Formel kann auch k"urzer definiert werden:

\starttypen
\chemie{2H_2,PLUS,O_2,GIVES,2H_2O}
\stoptypen

oder selbst:

\starttypen
\chemie{2H_2,+,O_2,->,2H_2O}
\stoptypen

Der erfahrene \TEX||Nutzer wird den erkennen, da"s sowohl das Pluszeichen als 
auch der Reaktionspfeil sich auf der Grundlinie befinden. Vergleichen Sie $+$ 
mit \chemie{+}. So wird daf"ur gesorgt, da"s unabh"angig von der Gr"o"se der
Formel das \chemie{+} und der \chemie{->} auf gleicher H"ohe liegen.

Neben \type{PLUS} und \type{GIVES} kann \type{EQUILIBRIUM} (\type{<->}) 
abgegeben werden, was einen Doppelpfeil \chemie{EQUILIBRIUM} liefert. 
Mit \type{MESOMERIC} oder \type{<>} bekommt man \chemie{MESOMERIC}.

Diese Formeln k"onnen auch im laufenden Text gesetzt werden. In diesem Fall
wird ein kleineres Layout ausgew"ahlt: \chemie{2H_2,+,O_2,->,2H_2O}. Es ist
auch m"oglich Bindungen anzuzeigen. Zum Beispiel ergibt 
\type{\chemie{H,SINGLE,CH,DOUBLE,HC,SINGLE,H}} die Formel
\chemie{H,SINGLE,CH,DOUBLE,HC,SINGLE,H}. Es geht auch k"urzer:
\type{\chemie{H,-,CH,--,HC,-,H}}. Eine Dreifachbindung erh"alt man mit
\type{TRIPLE} oder \type{---}: \chemie{HC,TRIPLE,CH}. 

Nun zur"uck zum Displaymodus. Das zweite Argument kann dazu verwendet werden,
um die Stoffe zu bezeichnen:

\startbuffer
\plaatsformule
  \startformule
    \chemie{2H_2}{Wasserstoff} \chemie{PLUS} \chemie{O_2}{Sauerstoff}
    \chemie{GIVES}{heftig} \chemie{2H_2O}{Wasser}
   \stopformule
\stopbuffer

\starttypen
$$
 \chemie{2H_2}{Wasserstoff} \chemie{PLUS} \chemie{O_2}{Sauerstoff}
 \chemie{GIVES}{heftig} \chemie{2H_2O}{Wasser}
$$
\stoptypen

Wir k"onnen dies "uber (oder unter) den Symbolen platzieren.

\haalbuffer

Es sind maximal drei Argumente m"oglich, wobei das letzte Argument nach unten
plaziert wird.

\startbuffer
\plaatsformule
  \startformule
    \chemie{H_2O}{liquid}{Wasser}
    \hbox{oder}
    \chemie{H_2O}{Wasser}
  \stopformule
\stopbuffer

\haalbuffer

Obrige Formel wurde so erstellt:

\starttypen
$$
 \chemie{H_2O}{liquid}{Wasser}
 \hbox{oder}
 \chemie{H_2O}{Wasser}
$$
\stoptypen

Die Gr"o"se der im Text gesetzten Formeln kann man mit dem Einstellungsbefehl
ge"andert werden.

\plaatstabel
{Einstellungen bei Textformeln.}
\starttabel[|lT|Tl|lT|]
\HL
\VL \bf Variable  \VL \bf Einstellungen  \VL \bf Vorgabe \VL\SR
\HL
\VL textgroesse   \VL klein mittel gross \VL gross       \VL\SR
\HL
\stoptabel

Der Befehl \type{\chemie{H,SINGLE,CH,DOUBLE,HC,SINGLE,H}} liefert
bei den Einstellungen \type{gross}, \type{mittel} und \type{klein} diese
Formeln:

\leavevmode\hbox
  {\stelchemiein[tekstformaat=groot]\chemie{H,-,CH,--,HC,-,H}\hskip2em
   \stelchemiein[tekstformaat=middel]\chemie{H,-,CH,--,HC,-,H}\hskip2em
   \stelchemiein[tekstformaat=klein]\chemie{H,-,CH,--,HC,-,H}}

\hoofdstuk{Indizies}

Tief- und hochgestellter Text wird, anders als es in Knuths {\TeX}Book
empfohlen ist, etwas tiefer gesetzt. So kann das auf Seite~179 des {\TeX}Books
gegebene Beispiel, das chemisch gesehen etwas merkw"urdig ist, mit
\type{\chemie{Fe_2^{+2}Cr_2O_4}} gesetzt werden. Der Befehl liefert 
\chemie{Fe_2^{+2}Cr_2O_4}, ohne Korrektur erhielte man $\rm Fe_2^{+2}Cr_2O_4$.

Der Platz der Indizies (tiefgestellter Text) wird per Standard der Variabel
\type{offset} entnommen: \type{HIGH} oder \type{LOW}. Diese Einstellungen
k"onnen mit den gleichnamigen Parametern "uberschrieben werden.

\startbuffer
\startchemie[kader=aan,optie=test,variant=2]
  \chemie[SIX,B,R13,HIGH,RZ1,LOW,RZ3][NH_3,NH_3]
\stopchemie
\stopbuffer

\startfiguurtekst
  [links][]
  {geen}
  {\haalbuffer}
\voorbeeld[vb:test]
\starttypen
\startchemie[rahmen=an,option=test,alternative=2]
  \chemie[SIX,B,R13,HIGH,RZ1,LOW,RZ3][NH_3,NH_3]
\stopchemie
\stoptypen
\stopfiguurtekst

Obwohl der Platz der Indizies f"ur jeden Substituenten so eingestellt werden 
kann, wird empfohlen, dies global mit \type{\stellechemieein} zu tun, damit 
die Formeln konsistent gesetzt werden.

Die Schl"usselw"orter \type{LOW} und \type{HIGH} k"onnen auch in 
Textformeln verwendet werden, so ergibt \type{\chemie{HIGH,H_2O}} 
\chemie{HIGH,H_2O} und \type{\chemie{LOW,H_2O}} \chemie{LOW,H_2O}.

\vervolgblad
  {Hintergrund}
  [links=1500,breedte=4000]
  [FIVE,ROT3,B125,+SB3,-SB4,ER35,SR4,CRZ35,Z4,RZ4,Z0]
  [O,O,N,CH_3]

\hoofdstuk{Installation}

Das Paket \PPCHTEX\ ist in erster Linie f"ur den Gebrauch innerhalb von
\CONTEXT\ entwickelt worden. Innerhalb dieses Makropakets kann \PPCHTEX\ mit
so geladen werden:\voetnoot{Die eigentlichen Makros stehen in der Datei
\type{ppchtex.tex}, die in \type{m-chemie.tex} (das~\type{m} steht hier f"ur
Modul) geladen werden.} 

\starttypen
\verwendemodule[pictex,chemie]
\stoptypen

Der Befehl l"adt zuerst die \PICTEX||Makros, so da"s \PPCHTEX\ wei"s welche
Ausgabe ge\-w"unscht ist. Die Chemiemakros sind somit nicht standardgem"a"s
verf"ugbar. 

Zudem kann das Paket auch in Kombination mit \PLAIN\ \TEX\ oder \LATEX\
verwendet werden. In letzterem Fall wird die Datei \type{m-chemie.sty}
verwendet; das Laden der \type{.sty}||Datei geschieht dann so: 

\starttypen
\documentstyle[m-pictex,m-chemie]{...}
\stoptypen

In \LATEX$2_{\textstyle\varepsilon}$ verwendet man:

\starttypen
\documentclass{...}
\usepackage{m-pictex}
\usepackage{m-chemie}
\stoptypen

Die Datei \type{m-pictex} sorgt daf"ur, da"s \PICTEX\ so effizient wie
m"oglich geladen wird. Das ist n"otig, weil \LATEX\ und auch die 
zus"atzlich geladenen Stildateien (\type{.sty}) ziemlich viele 
\type{\dimen}s reservieren. "Ubrigens ist als Folge der 
Benutzerfreundlichkeit auch eine gr"o"sere \TEX||Version n"otig.
Meistens wird dies bevorzugt.

\PPCHTEX\ ist in mehrern Sprachen verf"ugbar. Die verwendete Sprache h"angt
davon ab, wie es geladen wird:

\startregelcorrectie
\starttabel[|l|l|]
\NC niederl\"andisch \NC \type{m-che-nl} = \type{m-chemie.sty} \NC\FR
\NC englisch         \NC \type{m-che-en} = \type{m-chemic.sty} \NC\MR
\NC deutsch          \NC \type{m-che-de}                       \NC\LR
\stoptabel
\stopregelcorrectie

In der Datei \type{ppchtex.noc} ist zu sehen, wie au"serhalb von \CONTEXT\ 
die Kopplung mit dem Hauptmakropaket realisiert ist. Diese Datei l"adt zuerst 
einige der Systemmodule von \CONTEXT\ und anschlie"send das, was in dem 
jeweiligen Paket nicht verf"ugbar ist.

Die Distribution besteht aus den Definitionsdateien 
\type{ppchtex.tex} und \type{ppchtex.noc}, den Aufrufdateien 
\type{m-che-nl.tex}, \type{m-che-en.tex} und \type{m-che-de.tex}, sowie
die alternativen Aufrufdateien \type{m-chemie.tex} und \type{m-chemic.tex}.

\hoofdstuk{Erweiterungen}

Es ist den Benutzern von \PPCHTEX\ nat"urlich grunds"atzlich
gestattet die Makros, auch auf andere Weise weiter zu verwenden.
Einige Vorsicht ist dabei aber geboten, da die Makros noch weiter
entwickelt, optimiert und robuster gemacht werden. Einige Makros
scheinen zwar weitgehend fertig zu sein, aber der Anschein t"auscht.
Beispielsweise unterst"utzt die Makrogruppe \type{stelle..ein} 
Verschachtelung und verschiedene \ASCII||Layouts. Vergleichen
Sie beispielsweise:

\start
\steltypenin[optie=kleur]
\starttypen
\stellechemieein[groesse=klein]
\stoptypen
\stop

mit:

\start
\steltypenin[optie=kleur]
\starttypen
\stellechemieein
  [groesse=klein,
   format=500,
   textgroesse=gross]
\stoptypen
\stop

Die Einstellungen k"onnen in beliebiger Reihenfolge "ubergeben werden.
Nach M"oglichkeit werden Leerzeichen und Zeilenumbr"uche unterdr"uckt und
Fehlermeldungen angezeigt.

Da \PPCHTEX\ in erster Linie als Modul von \CONTEXT\ gedacht (und entwickelt)
worden ist, unterst"utzt das Paket mehrere Sprachen. Im Moment gibt es
ein niederl"andisches, ein englischsprachiges und ein deutschsprachiges
Interface.

Ein Blick auf die Datei \type{ppchtex.tex} zeigt, da"s beim Interpretieren
der Angaben zwischen den \type{[ ]} bei \type{\chemie} die \type{\processaction}||Makros
benutzt werden. Diese Makros sind, weil hier sowohl Verschachtelung als
auch ein "ubersichtliches Layout m"oglich ist, relativ langsam. Dem gegen"uber
steht der Vorteil, das \PPCHTEX||Makros recht intuitiv sind. In der Zwischenzeit
sind die Makros selbst und ihr Gebrauch weitgehend optimiert. Tests mit der
Stopuhr haben gezeigt, da"s weitere Optimierungen kaum zu einem Zeitvorteil
f"uhren, und es dabei dies zum Verlust von Flexibilit"at kommt.

Bei der Entwicklung von \PPCHTEX wurde die Weise, in der die Befehle
innerhalb der Produktionsumgebung \TEXEDIT\ wiedergegeben werden,
ber"ucksichtigt. In diesem Programm werden \TEX||Befehle und ||symbole
mit einer besonderen Funktion farbig dargestellt (wie hier"uber 
beispielhaft dargestellt).

\hoofdstuk{Schrift}

Die Makros verwenden den mathematischen Modus und somit \type{\textfont},
\type{\scriptfont} und \type{\scriptscriptfont}. Wenn es n"otig ist werden
auch die \type{\fontdimen}s~14, 16 und~17 von \type{\font2} angepa"st.
Dabei wird die aktuelle Flie"stextgr"o"se ber"ucksichtigt, die der $x$-H"ohe 
(\type{\fontdimen5}) entnommen wird.

"Anderungen in den \type{\fontdimen}s haben globalen Charakter, somit ist
Gruppierung sinnlos. Aus diesem Grund werden die Dimensionen immer gesetzt
und zur"uckgesetzt. Die derzeitige L"osung mag nicht die beste sein, aber
andere f"uhren zu Problemen mit skalierten Schriften.

Das Platzieren von Atomen und Molek"ulen (Text) kostet relativ viel
Zeit bei der Berechnung. In \CONTEXT\ h"angt die Geschwindigkeit auch von
der Komplexit"at des Makros \type{\rm} ab. Eine effizientere L"osung ist
m"oglich.

\hoofdstuk{Definitionen}

\PPCHTEX\ greift f"ur die Interaktivit"at auf \CONTEXT\ zur"uck. Dies bedeutet,
da"s das Interface mehrsprachig ist. Gegen"uber dem Vorteil, da"s jeder in
seiner Muttersprache mit dem Makropaket kommunizieren kann, steht der Nachteil,
da"s der gemeinschaftliche Gebrauch der Makros etwas umst"andlicher ist.

Zur Zeit sind viele der zugrunde liegende Befehle von \CONTEXT\
niederl"andisch. Das gleiche gilt f"ur die Parameter. Die zugrunde
liegenden Befehle von \PPCHTEX\ sind inzwischen englisch. Um 
beispielsweise Bibliotheken international gebrauchen zu k"onnen, m"ussen
diese mit den englischsprachigen Befehlen definiert werden, zum Beispiel:

\starttypen
\startchemical
  \chemical[SIX,B,C]
\stopchemical
\stoptypen

Einstellungen sind ein wenig umst"andlicher. Wir benutzen daf"ur 
Systemkonstanten und -variablen, die zur Zeit noch niederl"andisch sind.
Demn"achst werden diese auch auf englisch sein.
Gl"ucklicherweise sind die Konstanten und Variabeln gut zu erkennen,
so da"s das Anpassen nicht schwierig ist.

\starttypen
\setupchemical[\c!breedte=10cm,\c!hoogte=\v!passend]
\stoptypen


Parameter werden mit einem vorangestelltem \type{\c!} und Einstellungen mit
einem \type{\v!} angegeben. Dies funktioniert nur wenn \type{!} ein Zeichen 
ist. Deshalb werden diese Einstellungen in ein \type{\unprotect}||\type{\protect}||Paar
eingeschlossen. So zum Beispiel:

\starttypen
\unprotect

\setupchemical[\c!breedte=10cm,\c!hoogte=\v!passend]

\startchemical
  \chemical[SIX,B,C]
\stopchemical

\protect
\stoptypen

Mehr Informationen "uber das Interface und die Arbeit an diesem ist der 
Dokumentation zu den \CONTEXT||Modulen und den \type{mult}||Gruppen zu entnehmen.

\hoofdstuk{Farbe}

Innerhalb von \CONTEXT\ ist es m"oglich, Teile einer Struktur zu kolorieren.
Im  \in{Beispiel}[vb:kleur] ist sowohl der Substituent als auch dessen 
Bindung rot gef"arbt.

\startbuffer
\startchemie[assenstelsel=aan,kader=aan,breedte=5000]
  \chemie[SIX,B]
  \kleur[rood]{\chemie[SIX,R1,RZ1][COOH]}
\stopchemie
\stopbuffer

\startfiguurtekst
  [links][]
  {geen}
  {\haalbuffer}
\voorbeeld[vb:kleur]
\starttypen
\startchemie[achsen=an,rahmen=an,breite=5000]
  \chemie[SIX,B]
  \farbe[rot]{\chemie[SIX,R1,RZ1][COOH]}
\stopchemie
\stoptypen
\stopfiguurtekst

Zuerst mu"s man jedoch mit dem \CONTEXT||Befehl \type{\stellefarbeein[status=start]} den
 Farbmechanismus aktivieren.

\hoofdstuk[txt:cooh]{Interaktivit"at}

Wenn man \PPCHTEX\ zusammen mit \CONTEXT\ verwendet, wird das Erstellen 
interaktiver Texte unterst"utzt. Unter einem interaktiven Text verstehen
wir einen Text, in dem man nachschlagen kann, w"ahrend Teile des Textes
aktiv sein k"onnen. Das hei"st, da"s ein Klick auf ein Wort zum
Sprung zur einer Erkl"arung f"uhrt. So kann beispielsweise durch einen
Klick im Text auf \naar{\chemie{COOH}}[sub:cooh] zu einem Sprung zu dem
zugeh"origen Substituenten f"uhren. Umgekehrt k"onnte ein Klick auf den
Substituenten zu einem Sprung zur Erkl"arung jenes f"uhren.

\startbuffer
\startchemie[assenstelsel=aan,kader=aan,breedte=5000]
  \chemie[SIX,B]
  \chemie[sub:cooh][SIX,R1,RZ1][COOH]
\stopchemie
\stopbuffer

\startfiguurtekst
  [links][]
  {geen}
  {\haalbuffer}
\voorbeeld[vb:interactie]
\starttypen
\startchemie[achsen=an,rahmen=an,breite=5000]
  \chemie[SIX,B]
  \chemie[sub:cooh][SIX,R1,RZ1][COOH]
\stopchemie
\stoptypen
\stopfiguurtekst

Wie man sieht wird ein drittes Argument angegeben: der Verweis 
\type{[sub:cooh]}. Das hei"st, man kann von dem Text auf den Substituent
\naar{\chemie{COOH}}[sub:cooh] verweisen:

\starttypen
... text ... \zu{\chemie{COOH}}[sub:cooh] ... text ...
\stoptypen

\type{\zu} ist hierbei ein \CONTEXT||Befehl. Das Klicken auf einen so 
markierten Text positioniert den Cursor "uber dem Substituenten der
Graphik.

Dies funktioniert auch umgekehrt. Wir k"onnen von einer Struktur auf einen
Text verweisen. Der Verweis mu"s nat"urlich bestehen.

Klicken auf \chemie{COOH} in der Figur bewirkt ein Springen zum zugeh"origen
Text, der zum Beispiel so markiert ist:

\starttypen
\absatz[txt:cooh]{Substituenten}

... text ... \chemie{COOH} ... text ... 
\stoptypen 

Eine Kombination ist auch m"oglich. In diesem Fall sollte man auf die erste 
mit \type{\chemie} verweisen und auf die n"achste mit \type{\zurchemie}. 

Die Verkn"upfung mit dem Interaktionsmechanismus geschieht mit diesen 
Makros:

\starttypen
\localgotochemical   {verweis} {text}
\localthisischemical {verweis}
\stoptypen

F"ur eine weitere Erl"auterung verweisen wir auf den Quelltext. Diese
Kopplung kann auch f"ur andere Zwecke benutzt werden. Dann sollte man jedoch
einige Erfahrung mit \TEX\ haben.

\vervolgblad
  {"Ubersicht}
  [links=1250,breedte=4000,onder=1500,hoogte=4000]
  [SIX,B,C,MOV1,B,ER6,CRZ6]
  [O]

\input man-comm.tex

\stoptekst
